\documentclass[prd,twocolumn,floatfix,nofootinbib]{revtex4}
\usepackage{amssymb}
\usepackage{amsmath}
\usepackage{graphicx,color,dcolumn,booktabs,bm}
\usepackage{longtable,lscape}
\usepackage{txfonts}
\usepackage{overpic}
\usepackage{indentfirst}
\usepackage{cases}
\usepackage{multirow}
\usepackage{epstopdf}
\usepackage{ulem}
\usepackage[colorlinks,
            citecolor=blue,
            anchorcolor=red,
            menucolor=red,
            linkcolor=red,
            filecolor=red,
            runcolor=red,
            urlcolor=blue,
            frenchlinks=false]{hyperref}
%\newtheorem{theorem}{Theorem}
%\newtheorem{remark}[theorem]{Remark}
%\newtheorem{solution}[theorem]{Solution}
%\newtheorem{summary}[theorem]{Summary}
%\newenvironment{proof}[1][Proof]{\noindent\textbf{#1.} }{\ \rule{0.5em}{0.5em}}
%\usepackage{tikz}
\setcounter{MaxMatrixCols}{10}
\def\lrpartial{\buildrel\leftrightarrow\over\partial}
\graphicspath{{Figures/}} %
\allowdisplaybreaks

\linespread{1.1}
\begin{document}
%\title{Ground-state baryons and tetraquarks with one and two heavy quarks: Masses, magnetic moments}
\title{Masses and magnetic moments of open heavy hadrons in MIT bag model:  singly and doubly heavy baryons and tetraquarks}

\author{Wen-Xuan Zhang$^{1}$}
\author{Hao Xu$^{1,2}$}
\author{Duojie Jia$^{1,2}$
\footnote{Corresponding author}}
\email{jiadj@nwnu.edu.cn}
\author{XueQian Li$^{3}$}

\affiliation{ $^1$Institute of Theoretical Physics, College of Physics and
Electronic Engineering, Northwest Normal University, Lanzhou 730070, China \\
$^2$Lanzhou Center for Theoretical Physics, Lanzhou University, Lanzhou, 730000,China \\
$^3$School of Physics Science, Nankai University, Tianjin 300071, China \\ }


\begin{abstract}
We compute ground-state masses and magnetic moments of heavy baryons with strange number $S\leq1$ and tetraquarks with one and two open heavy flavors in MIT bag model. In addition to the standard bag model parameters, we use masses of heavy mesons with $J^{P}=1^{-}$ to extract the effective masses of charm and bottom quarks, and add binding energies assumed to exist among heavy quarks($c$ and $b$) and strange quarks in the literatures. Notable flavor-dependence is seen in further predictions of the bag radius and charge radii for all these open heavy baryons and tetraquarks, including the doubly charmed baryons $\Xi_{cc}^{\ast}$ and their iso-partners, as well as the strange isosinglet tetraquark $ud\bar{s}\bar{c}$ with $J^{P}=0^{+}$ and mass  $M\left(ud\bar{s}\bar{c}\right)=2.934$ GeV. Our numerical variation indicates that chromomagnetic mixing in $\Xi_{c}$ , $\Xi_{b}$, $\Xi_{bc}$, $\Omega_{bc}$  with $J^{P}={\frac{1}{2}}^{+}$ and among is sizable, compared to perturbative treatment for chromomagnetic interaction.

{\normalsize PACS number(s):12.39Jh, 12.40.Yx, 12.40.Nn}
{\normalsize Key Words: Heavy tetraquark, Spectroscopy, Quantum number, Magnetic moment}
\end{abstract}
\maketitle
\date{\today}


\maketitle\section{Introduction}
In 2017, the LHCb Collaboration at CERN discovered the doubly charmed baryon $\Xi_{cc}^{++}$ with $J^{P}={1/2}^{+}$ and the measured mass $3621.40\pm 0.78$ MeV \cite{Aaij:2017ueg}. This observation provides important experimental information about the strength of the interaction between two heavy quarks and makes it possible to resolve the question whether the tetraquarks $QQ\bar{q}\bar{q}$ with two open heavy quarks $Q$ and two light antiquarks $\bar{q}$ are stable or unstable against decay into two $Q\bar{q}$ mesons \cite{KR:prl17,EichtenQ:prl17,LuoCL:Epj17}. More recently, LHCb Collaboration reported the first exotic state $X_{0}(2900)$ with open heavy flavor and mass $2866\pm 7$ MeV \cite{LHCb:2020kd}. It has been interpreted as an isosinglet tetraquark $cs\bar{u}\bar{d}$ in Ref.~\cite{Karliner:2020vsi}. These foundings aroused great interest in doubly heavy baryons and tetraquarks.

Various approaches have been employed to explore doubly heavy baryons and tetraquarks in the literatures, including potential quark models and bag model \cite{Fleck:1989mb,Ebert:2004ck,Bernotas:2012nz,He:2004px,Karliner:2014gca,Liu:2018euh} for doubly heavy baryons. For the tetraquarks, see Refs. \cite{AliLS:ppn2017,Liu:2019zoy} for recent reviews. Information like mass spectrum and other properties for heavy hadrons can be important in identifying and/or finding them eventually. For instance, a $\Xi_{cc}$ candidate with a mass about $100$ MeV lighter than the quark model prediciton about $3620$ MeV\cite{Ebert:2004ck,Karliner:2014gca} was observed previously in 2002 by the SELEX Collaboration at Fermilab\cite{Moinester:2002uw} and has not been confirmed by other experiments.

In this work, we employ MIT bag model to systematically study masses and magnetic moments of singly and doubly heavy baryons with strange number $S\leq1$ and tetraquarks with one and two open heavy flavors. By the way, the masses and other properties of the ground state light mesons, heavy mesons are also calculated. For possible chromomagnetic mixing in $\Xi_{c}$ , $\Xi_{b}$, $\Xi_{bc}$, $\Omega_{bc}$  with $J^{P}={\frac{1}{2}}^{+}$ and among, we computed via numerical variation the splitted magnetic moments and charge radii of them, compared to perturbative treatment of chromomagnetic interaction.

In the past few decades, bag model have been applied to address mass spectrum and other properties of the doubly heavy baryons \cite{Fleck:1989mb,He:2004px, Bernotas:2012nz} and multiquark hadrons, including light multiquark baryons with five and seven nonstrange quarks \cite{Strottman:1979qu}. A slowly running strong coupling $\alpha_{s}$ was applied to evaluate the mass spectra of doubly heavy baryons\cite{He:2004px} and a Coulomb-like interaction between heavy quarks in a bag is derived in Refs.\cite{Haxton:1980mc,Aerts:1980rf}.

It is well known that bag model\cite{DeGrandJ:1975D}, simple though it is, embodies two primary features of quantum chromodynamics(QCD): asymptotic freedom at short distance and confinement at long distance. We find that the MIT bag model, when equiped with a slowly running strong coupling $\alpha_{s}$ and binding energies which are assumed to exist among heavy quarks($c$ and $b$) and strange quarks in the literatures, reconciles the computation of the heavy hadrons with strange number $S\leq1$ with the known results for the light hadron. The mass errors of our calculation reduce to about  MeV for heavy mesons and to about $40$ MeV for heavy baryons, from which we proceed to predict the masses and other properties of the tetraquarks with one and two open heavy quarks.

This paper will be organized as follows. In Sec.~\ref{Method}, we review some basic relations of MIT bag model, including chromomagnetic interaction (CMI) among the quarks in bag. In Sec.~\ref{Parameters}, a systematic numerical calculation is performed for the established light hadrons and heavy baryons, with the optimal set of parameters obtained and the results for the mass spectra and other properties reproduced.
In Sec.~\ref{Hadrons}, we present a detailed predictions for masses and other propeties for doubly heavy baryons and doubly (open) heavy tetraquarks. The paper ends with summary and conclusions in Sec.~\ref{Discussions}.


\section{Method for MIT bag model and CMI}\label{Method}
\subsection{Mass Formula}
Treated hadron as a sphere bag, MIT bag model provides an approach to estimate masses and other properties of hadrons in their ground-states \cite{DeGrandJ:1975D,DeGrand:1975cf}, in which the chromomagnetic interaction derived from the energy of a sphere-like gluon field, interacted with quark fields in bag\cite{DeGrand:1975cf}. The mass formula of hadron in MIT bag model is\cite{DeGrand:1975cf},
\begin{equation}\label{MBm}
    M\left(R\right)=\sum_{i=n,s,c,b} n_{i}\omega_{i}+\frac{4}{3}\pi R^{3}B-\frac{Z_{0}}{R}+\Delta M
\end{equation}
\begin{equation}\label{freq}
    \omega_{i}=\left(m_{i}^{2}+\frac{x_{i}^{2}}{R^2}\right)^{1/2}
\end{equation}
where the first term the kinematic energy of all quarks, the second is volume energy, the third is the zero-point-energy (ZPE) and $\Delta M$ is interaction energy among quarks in bag, on which we will address in this work. Here in Eq. (\ref{MBm}), $n_{i}$ is number of quark or antiquark in bag with mass $m_{i}$ and flavor $i$, where $i$ can be $n=u,d$, $s$, $c$ and $b$, $x_{i}=x_{i}(mR)$ is dimensionless dynamical parameters, $R$ the bag radius to be determined variationally, $B$ the bag constant and $Z_{0}$ is zero-point-energy (ZPE) constant. $x_{i}$ and $R$ are related by a transcendental equation
\begin{equation}\label{transc}
    \tan x_{i}=\frac{x_{i}}{1-m_{i}R-\left(m_{i}^{2}R^{2}+x_{i}^{2}\right)^{1/2}}
\end{equation}

For the charm and bottom quarks($i=c,b$), $\omega_{c,b}$ is approximated by $\omega_{c,b}\simeq m_{c,b}$. The interaction energy $\Delta M$ composed of two energy terms:
(1) The short-ranged chromoelectric interaction energy, due mainly to the perturbative gluon exchange between quarks. Owing to its smallness for the relativistic light quarks involving the nonstrange quark $n(=u,d)$, this energy, scales mainly as $-\sum\alpha_{s}/r_{ij}$, becomes sizable when both of two quarks $i$ and $j$ are massive and moving nonrelativistically. In present work, we treat this energy as sum of binding energies among heavy quarks and between heavy quark $Q$ and strange quark $s$\cite{Karliner:2014gca,Karliner:17N,Karliner:18D}.
The net effect for this chromoelectric interaction amounts to introduction of five binding energies $B_{QQ\prime}$ and $B_{Qs}$: $B_{cs}$, $B_{cc}$, $B_{bs}$, $B_{bb}$ and $B_{bc}$, which can be extracted from heavy mesons for the pair configuration $\boldsymbol{\bar{3}}_{c}$ and scaled to other color configurations.

(2) The chromomagnetic interaction energy, due to short-ranged perturbative gluon exchange between quarks (antiquark) $i$ and $j$,
\begin{equation}\label{CMI}
    H_{CMI}=-\sum_{i<j} \left(\boldsymbol{\lambda_{i}}\cdot \boldsymbol{\lambda_{j}}\right) \left(\boldsymbol{\sigma_{i}}\cdot \boldsymbol{\sigma_{j}}\right)C_{ij}
\end{equation}
where $\boldsymbol{\lambda}_{i}$ are Gell-Mann matrices and  $\boldsymbol{\sigma}_{i}$ the Pauli matrices, $C_{ij}$ is the CMI parameter. This parameter is given, in MIT bag model, by
\begin{equation}\label{Cij}
    C_{ij}=3\frac{\alpha_{s}\left(R\right)}{R^{3}}\bar{\mu}_{i}\bar{\mu}_{j}I_{ij}
\end{equation}
\begin{equation}\label{muBari}
   \bar{\mu}_{i}=\frac{R}{6} \frac{4\alpha_{i}+2\lambda_{i}-3}{2\alpha_{i}\left(\alpha_{i}-1\right)+\lambda_{i}}
\end{equation}
\begin{equation}\label{Iij}
    I_{ij}=1+2\int_{0}^{R} \frac{{\rm d}r}{r^{4}} \bar{\mu}_{i}\bar{\mu}_{j}=
    1+F(x_{i} x_{j})
\end{equation}
with 
\begin{equation}\label{F}
    F(x_{i} x_{j})= ???
\end{equation}
where $\alpha_{i}\equiv\omega_{i}R$, $\lambda_{i}\equiv m_{i}R$, $\alpha_{s}\left(R\right)$ is running strong coupling constant, $\bar{\mu}_{i}$ is reduced magnetic moment without quark charge,
and $I_{ij}$ takes a rational function of $x_{i}$ and $x_{j}$ Ref.~\cite{DeGrand:1975cf}. 

In some of bag models \cite{Haxton:1980mc,Aerts:1980rf,Carlson:1982er,He:2004px,Fleck:1989mb}, the parameter $\alpha_{s}$ takes, inspired by QCD renormalization, a semiempirical form like
\begin{equation}\label{alphaS}
    \alpha_{s}(R)=\frac{2\pi n}{9} \frac{1}{ln\left[\gamma+{\left(R\Lambda_{QCD}\right)}^{-n}\right]}
\end{equation}
where exponent $n=1$ or 2, with 1 prefered for suppressing the value of $\alpha_{s}$, $\Lambda_{QCD}$ is the QCD scale ranging from $0.2$ GeV to $0.5$ GeV, and $\gamma$ is prefactor used to avoid divergence. We take $\Lambda_{QCD}=0.281$ GeV to set
\begin{equation}\label{alphaS-mine}
    \alpha_{s}(R)=\frac{0.296}{ln\left[1+{\left(0.281R\right)}^{-1}\right]}
\end{equation}
which is plotted in FIG.~\ref{fig:alphaS}. Among four running behaviors of $\alpha_{s}$ in the plot, the solid line provides notable variation and simultaneously gives a relative lower value of $\alpha_{s}$ no more than $0.6$ at $R=5$ GeV$^{-1}$. In the case of light hadrons, mass spectra evaluated by Eq.~(\ref{alphaS-mine}) are listed in Table \ref{tab:lighthadron}, compared to that given by original MIT bag model.

\begin{figure}[!ht]
    \centering
    \begin{tabular}{c}
        \includegraphics[width=0.4\textwidth]{Figures/AlphaS.eps}
    \end{tabular}
    \caption{Four running behaviors of strong coupling constant.
    Bag radius $R$ takes range from $3\,{\rm GeV}^{-1}$ to $6\,{\rm GeV}^{-1}$, and the standard radius is $5\,{\rm GeV}^{-1}$ ($\approx 1\,{\rm fm}$) for checking.
    The solid line represents our result (\ref{alphaS-mine}). Dashed line draws Eq.~(\ref{alphaS}) when $\gamma=2.847$ and $\Lambda_{QCD}=0.281\,{\rm GeV}$.
    Dotted line draws Eq.~(\ref{alphaS}) when $\gamma=1$ and $\Lambda_{QCD}=0.281\,{\rm GeV}$. The dotdashed line indicates Eq.~(\ref{alphaS-2004}).}
    \label{fig:alphaS}
\end{figure}

Given the parameters, the quark mass $m_{i}$, bag constant $B$, the ZPE coefficient $Z_{0}$ and strong coupling constant $\alpha_{s}(R)$ depending on bag radius $R$, one can determine, via numerical variational method, the respective bag radius $R$ for each hadron, and thereby compute the ground-state masses and other static properties, like magnetic moments, the charge radius, from mesons to tetraquarks, by variational method. 

Whenever the chromomagnetic mixing occured, e.g., for the degenerated states in  $\Xi_{Q}$ , $\Xi_{bc}$, $\Omega_{bc}$  with $J^{P}={\frac{1}{2}}^{+}$ and among $ns\bar{c}\bar{}$, we firstly diagonalize the mass in Eq. (\ref{MBm}) and then solve the bag radius $R$ variationally. Meanwhile, a first-order perturbative  computation is done for them for comparison treating CMI as a pertubation. 


\subsection{Chromomagnetic Structure}
To evaluate the mass splitting due to the chromomagnetic interaction in Eq.~(\ref{CMI}), in which $\boldsymbol{\lambda}_{i}$ should be replaced by $-\boldsymbol{\lambda}_{i}$ for an antiquark, the CMI matrices have to be computed for given hadron with certain $J^{P}$ quantum number so that its diagonalizing gives the respective mass splittings.  \cite{Luo:2017eub}, Here, we list all respective flavor-spin-color wave functions for meson, baryon and tetraquark systems considered in this work and present relevant formulas for the CMI matrix elements for the hadrons with given chromomagnetic structures and without chromomagnetic mixing.

A ground-state meson with color function $\chi_{0}^{M}$)  
\begin{equation}\label{colorM}
    \phi^{M}=\left| q_{1}\bar{q}_{2} \right\rangle
\end{equation}
can have two spin states(one of vector like, $\chi_{1}^{M}$ and other of scalor like):
\begin{equation}\label{spinM}
    \chi_{1}^{M}={\left| q_{1}\bar{q}_{2} \right\rangle}_{1} \quad \chi_{2}^{M}={\left| q_{1}\bar{q}_{2} \right\rangle}_{0}
\end{equation}
where subscript $J$ outside the bracket denotes the spin of hadron. The color-spin wavefunctions of a meson with spin $J=0,1$ are
\begin{equation}\label{colorspinM}
    \phi^{M}\chi_{1}^{M}={\left| q_{1}\bar{q}_{2} \right\rangle}_{1} \quad \phi^{M}\chi_{2}^{M}={\left| q_{1}\bar{q}_{2} \right\rangle}_{0}
\end{equation}

A baryon with color function  
\begin{equation}\label{colorB}
    \phi^{B}=\left| {\left(q_{1}q_{2}\right)}^{\bar{3}} q_{3} \right\rangle
\end{equation}
can have three spin states
\begin{equation}\label{spinB}
    \begin{aligned}
        \chi_{1}^{B}={\left| {\left(q_{1}q_{2}\right)}_{1} q_{3} \right\rangle}_{3/2} \\
        \chi_{2}^{B}={\left| {\left(q_{1}q_{2}\right)}_{1} q_{3} \right\rangle}_{1/2} \\
        \chi_{3}^{B}={\left| {\left(q_{1}q_{2}\right)}_{0} q_{3} \right\rangle}_{1/2}
    \end{aligned}
\end{equation}
where $\left(q_{1}q_{2}\right)$ stands for a diquark in color $\boldsymbol{\bar{3}}_{c}$ with spin $J=1$ or $0$.

To write the flavor-spin-color function for a hadron, one has to consider the flavor symmetry. For a symmetric flavor wave function of $\left(q_{1}q_{2}\right)$, with isospin $I=1$ or flavors identical, we use a symbol $\delta_{12}^{S}=1$. For an asymmetric flavor wave function, with $I=0$
there will be a symbol $\delta_{12}^{A}=1$. If two quarks $q_{1}$ and $q_{2}$ carry different flavors beyond SU(2) isospin symmetry, the wave function can be either symmetric or asymmetric, one can use $\delta_{12}^{S}=\delta_{12}^{A}=1$.
By Pauli principle, we get then three flavor-spin-color wave functions for a baryon
\begin{equation}\label{colorspinB}
    \begin{aligned}
        \phi^{B}\chi_{1}^{B}={\left| {\left(q_{1}q_{2}\right)}_{1}^{\bar{3}} q_{3} \right\rangle}_{3/2} \delta_{12}^{S} \\
        \phi^{B}\chi_{2}^{B}={\left| {\left(q_{1}q_{2}\right)}_{1}^{\bar{3}} q_{3} \right\rangle}_{1/2} \delta_{12}^{S} \\
        \phi^{B}\chi_{3}^{B}={\left| {\left(q_{1}q_{2}\right)}_{0}^{\bar{3}} q_{3} \right\rangle}_{1/2} \delta_{12}^{A}
    \end{aligned}
\end{equation}

A possible state mixing can enter due to chromomagnetic interaction. 
For example, for $\Xi_{bc}$ and $\Xi_{bc}^{\prime}$ with flavor structure $(bc)s$ choosing $\delta_{12}^{S}=\delta_{12}^{A}=1$ can not distinguish $\phi^{B}\chi_{2}^{B}$ and $\phi^{B}\chi_{3}^{B}$ states by their $J^{P}$ quantum number, and thus the physical state must be a mixing of $\phi^{B}\chi_{2}^{B}$ and $\phi^{B}\chi_{3}^{B}$. See more details about colormagnetic mixing in Sec.~\ref{Parameters}.

A tetraquark with color structure $\boldsymbol{6}\otimes\boldsymbol{\bar{6}}$ or $\boldsymbol{\bar{3}}\otimes\boldsymbol{3}$ 
\begin{equation}\label{colorT}
    \phi_{1}^{T}=\left| {\left(q_{1}q_{2}\right)}^{6} {\left(\bar{q}_{3}\bar{q}_{4}\right)}^{\bar{6}} \right\rangle \quad
    \phi_{2}^{T}=\left| {\left(q_{1}q_{2}\right)}^{\bar{3}} {\left(\bar{q}_{3}\bar{q}_{4}\right)}^{3} \right\rangle
\end{equation}
can have six spin states
\begin{equation}\label{spinT}
    \begin{aligned}
        \chi_{1}^{T}={\left| {\left(q_{1}q_{2}\right)}_{1} {\left(\bar{q}_{3}\bar{q}_{4}\right)}_{1} \right\rangle}_{2} \quad
        \chi_{2}^{T}={\left| {\left(q_{1}q_{2}\right)}_{1} {\left(\bar{q}_{3}\bar{q}_{4}\right)}_{1} \right\rangle}_{1} \\
        \chi_{3}^{T}={\left| {\left(q_{1}q_{2}\right)}_{1} {\left(\bar{q}_{3}\bar{q}_{4}\right)}_{1} \right\rangle}_{0} \quad
        \chi_{4}^{T}={\left| {\left(q_{1}q_{2}\right)}_{1} {\left(\bar{q}_{3}\bar{q}_{4}\right)}_{0} \right\rangle}_{1} \\
        \chi_{5}^{T}={\left| {\left(q_{1}q_{2}\right)}_{0} {\left(\bar{q}_{3}\bar{q}_{4}\right)}_{1} \right\rangle}_{1} \quad
        \chi_{6}^{T}={\left| {\left(q_{1}q_{2}\right)}_{0} {\left(\bar{q}_{3}\bar{q}_{4}\right)}_{0} \right\rangle}_{0}
    \end{aligned}
\end{equation}
which gives twelve base wave functions
\begin{equation}\label{colorspinT}
    \begin{aligned}
        \phi_{1}^{T}\chi_{1}^{T}={\left| {\left(q_{1}q_{2}\right)}_{1}^{6} {\left(\bar{q}_{3}\bar{q}_{4}\right)}_{1}^{\bar{6}} \right\rangle}_{2} \delta_{12}^{A}\delta_{34}^{A} \\
        \phi_{2}^{T}\chi_{1}^{T}={\left| {\left(q_{1}q_{2}\right)}_{1}^{\bar{3}} {\left(\bar{q}_{3}\bar{q}_{4}\right)}_{1}^{3} \right\rangle}_{2} \delta_{12}^{S}\delta_{34}^{S} \\
        \phi_{1}^{T}\chi_{2}^{T}={\left| {\left(q_{1}q_{2}\right)}_{1}^{6} {\left(\bar{q}_{3}\bar{q}_{4}\right)}_{1}^{\bar{6}} \right\rangle}_{1} \delta_{12}^{A}\delta_{34}^{A} \\
        \phi_{2}^{T}\chi_{2}^{T}={\left| {\left(q_{1}q_{2}\right)}_{1}^{\bar{3}} {\left(\bar{q}_{3}\bar{q}_{4}\right)}_{1}^{3} \right\rangle}_{1} \delta_{12}^{S}\delta_{34}^{S} \\
        \phi_{1}^{T}\chi_{3}^{T}={\left| {\left(q_{1}q_{2}\right)}_{1}^{6} {\left(\bar{q}_{3}\bar{q}_{4}\right)}_{1}^{\bar{6}} \right\rangle}_{0} \delta_{12}^{A}\delta_{34}^{A} \\
        \phi_{2}^{T}\chi_{3}^{T}={\left| {\left(q_{1}q_{2}\right)}_{1}^{\bar{3}} {\left(\bar{q}_{3}\bar{q}_{4}\right)}_{1}^{3} \right\rangle}_{0} \delta_{12}^{S}\delta_{34}^{S} \\
        \phi_{1}^{T}\chi_{4}^{T}={\left| {\left(q_{1}q_{2}\right)}_{1}^{6} {\left(\bar{q}_{3}\bar{q}_{4}\right)}_{0}^{\bar{6}} \right\rangle}_{1} \delta_{12}^{A}\delta_{34}^{S} \\
        \phi_{2}^{T}\chi_{4}^{T}={\left| {\left(q_{1}q_{2}\right)}_{1}^{\bar{3}} {\left(\bar{q}_{3}\bar{q}_{4}\right)}_{0}^{3} \right\rangle}_{1} \delta_{12}^{S}\delta_{34}^{A} \\
        \phi_{1}^{T}\chi_{5}^{T}={\left| {\left(q_{1}q_{2}\right)}_{0}^{6} {\left(\bar{q}_{3}\bar{q}_{4}\right)}_{1}^{\bar{6}} \right\rangle}_{1} \delta_{12}^{S}\delta_{34}^{A} \\
        \phi_{2}^{T}\chi_{5}^{T}={\left| {\left(q_{1}q_{2}\right)}_{0}^{\bar{3}} {\left(\bar{q}_{3}\bar{q}_{4}\right)}_{1}^{3} \right\rangle}_{1} \delta_{12}^{A}\delta_{34}^{S} \\
        \phi_{1}^{T}\chi_{6}^{T}={\left| {\left(q_{1}q_{2}\right)}_{0}^{6} {\left(\bar{q}_{3}\bar{q}_{4}\right)}_{0}^{\bar{6}} \right\rangle}_{0} \delta_{12}^{S}\delta_{34}^{S} \\
        \phi_{2}^{T}\chi_{6}^{T}={\left| {\left(q_{1}q_{2}\right)}_{0}^{\bar{3}} {\left(\bar{q}_{3}\bar{q}_{4}\right)}_{0}^{3} \right\rangle}_{0} \delta_{12}^{A}\delta_{34}^{A}
    \end{aligned}
\end{equation}
For all relevant color wave function see Appendix A and for relevant spin wave functions see Appendix B. Thus, one can calculate the color and spin factors in Eq.~(\ref{CMI}) with the help of the following relations
\begin{equation}\label{colorfc}
    {\left\langle \boldsymbol{\lambda_{i}}\cdot \boldsymbol{\lambda_{j}} \right\rangle}_{nm}=
    \sum_{\alpha=1}^{8} Tr\left(c_{in}^{\dagger}\lambda^{\alpha}c_{im}\right) Tr\left(c_{jn}^{\dagger}\lambda^{\alpha}c_{jm}\right)
\end{equation}
\begin{equation}\label{spinfc}
    {\left\langle \boldsymbol{\sigma_{i}}\cdot \boldsymbol{\sigma_{j}} \right\rangle}_{xy}=
    \sum_{\alpha=1}^{3} Tr\left(\chi_{ix}^{\dagger}\sigma^{\alpha}\chi_{iy}\right) Tr\left(\chi_{jx}^{\dagger}\sigma^{\alpha}\chi_{jy}\right)
\end{equation}
where $n$, $m$ and $x$, $y$ indicate the specific color and spin states respectively, $i$, $j$ the indexes of quarks or antiquarks, and the functions $c$ and $\chi$ are the respective basis vectors in the color and spin spaces.
A set of non-mixing hadronic states and its respective CMI are listed in Table \ref{tab:CMI}. 

Now, we are in the position to evaluate the matrix formula of the  the CMI energy (\ref{CMI}) and diagonalize it and thereby to minimize the obtained mass formula, after adding the bag energy and the binding energy (for heavy quark pair and for a heavy quark and a strange quark), to determine the dynamical parameters $x_{i}$ and $R$ and thereby obtain the full bag wavefunction for a hadron.  

\renewcommand{\tabcolsep}{0.49cm}
\renewcommand{\arraystretch}{1.5}
\begin{table*}[!htbp]
    \caption{CMI for the non-mixing hadrons with given color-spin structures. $C_{ij}$ follows Eq.~(\ref{Cij}) with subscripts which do not differ quark and antiquark.}
    \label{tab:CMI}
    \begin{tabular}{cccccc}
        \hline\hline
        {\rm State}     &{\rm Structure}            &{\rm CMI}      &{\rm State}        &{\rm Structure}            &{\rm CMI} \\ \hline
        $\pi$           &$\phi^{M}\chi_{2}^{M}$     &$-16C_{nn}$    &$\omega$           &$\phi^{M}\chi_{1}^{M}$     &$\frac{16}{3}C_{nn}$ \\
        $K$             &$\phi^{M}\chi_{2}^{M}$     &$-16C_{sn}$    &$K^{\ast}$         &$\phi^{M}\chi_{1}^{M}$     &$\frac{16}{3}C_{sn}$ \\
                        &                           &               &$\phi$             &$\phi^{M}\chi_{1}^{M}$     &$\frac{16}{3}C_{ss}$ \\
        $D$             &$\phi^{M}\chi_{2}^{M}$     &$-16C_{cn}$    &$D^{\ast}$         &$\phi^{M}\chi_{1}^{M}$     &$\frac{16}{3}C_{cn}$ \\
        $D_{s}$         &$\phi^{M}\chi_{2}^{M}$     &$-16C_{cs}$    &$D_{s}^{\ast}$     &$\phi^{M}\chi_{1}^{M}$     &$\frac{16}{3}C_{cs}$ \\
        $\eta_{c}$      &$\phi^{M}\chi_{2}^{M}$     &$-16C_{cc}$    &$J/\psi$           &$\phi^{M}\chi_{1}^{M}$     &$\frac{16}{3}C_{cc}$ \\

        $N$             &$\phi^{B}\chi_{3}^{B}$     &$-8C_{nn}$                                 &$\Delta$               &$\phi^{B}\chi_{1}^{B}$     &$8C_{nn}$ \\
        $\Lambda$       &$\phi^{B}\chi_{3}^{B}$     &$-8C_{nn}$                                 &                       &                           & \\
        $\Sigma$        &$\phi^{B}\chi_{2}^{B}$     &$\frac{8}{3}C_{nn}-\frac{32}{3}C_{sn}$     &$\Sigma^{\ast}$        &$\phi^{B}\chi_{1}^{B}$     &$\frac{8}{3}C_{nn}+\frac{16}{3}C_{sn}$ \\
        $\Lambda_{c}$   &$\phi^{B}\chi_{3}^{B}$     &$-8C_{nn}$                                 &                       &                           & \\
        $\Sigma_{c}$    &$\phi^{B}\chi_{2}^{B}$     &$\frac{8}{3}C_{nn}-\frac{32}{3}C_{cn}$     &$\Sigma_{c}^{\ast}$    &$\phi^{B}\chi_{1}^{B}$     &$\frac{8}{3}C_{nn}+\frac{16}{3}C_{cn}$ \\
        $\Xi$           &$\phi^{B}\chi_{2}^{B}$     &$\frac{8}{3}C_{ss}-\frac{32}{3}C_{sn}$     &$\Xi^{\ast}$           &$\phi^{B}\chi_{1}^{B}$     &$\frac{8}{3}C_{ss}+\frac{16}{3}C_{sn}$ \\
        $\Xi_{c}$       &$\phi^{B}\chi_{3}^{B}$     &$-8C_{sn}$                                 &$\Xi_{c}^{\ast}$       &$\phi^{B}\chi_{1}^{B}$     &$\frac{8}{3}C_{sn}+\frac{8}{3}C_{cn}+\frac{8}{3}C_{cs}$ \\
                        &                           &                                           &$\Omega$               &$\phi^{B}\chi_{1}^{B}$     &$8C_{ss}$ \\

        ${\left(nn\bar{c}\bar{c}\right)}^{I=1}$     &$\phi_{2}^{T}\chi_{2}^{T}$     &$\frac{8}{3}C_{nn}-\frac{16}{3}C_{cn}+\frac{8}{3}C_{cc}$
        &${\left(nn\bar{c}\bar{c}\right)}^{I=1}$    &$\phi_{2}^{T}\chi_{1}^{T}$     &$\frac{8}{3}C_{nn}+\frac{16}{3}C_{cn}+\frac{8}{3}C_{cc}$ \\
        ${\left(nn\bar{s}\bar{c}\right)}^{I=0}$     &$\phi_{1}^{T}\chi_{1}^{T}$     &$-\frac{4}{3}C_{nn}+\frac{20}{3}C_{sn}+\frac{20}{3}C_{cn}-\frac{4}{3}C_{cs}$
        &${\left(nn\bar{s}\bar{c}\right)}^{I=1}$    &$\phi_{2}^{T}\chi_{1}^{T}$     &$\frac{8}{3}C_{nn}+\frac{8}{3}C_{sn}+\frac{8}{3}C_{cn}+\frac{8}{3}C_{cs}$ \\
        \hline\hline
    \end{tabular}
\end{table*}

\renewcommand{\tabcolsep}{0.7cm}
\renewcommand{\arraystretch}{1.5}
\begin{table*}[!htbp]
    \caption{Sum rule for magnetic moment of given spin state or a mixing system.}
    \label{tab:musum}
    \begin{tabular}{lcc}
        \hline\hline
        {\rm Hadron}        &$\psi_{spin}$      &$\mu$ \\ \hline
        $\left(q_{1}{\bar{q}}_{2}\right)$
                            &$\chi_{1}^{M}$     &$\mu_{1}+\mu_{2}$ \\
                            &$\chi_{2}^{M}$     &0 \\
        $\left(q_{1}q_{2}\right)q_{3}$
                            &$\chi_{1}^{B}$     &$\mu_{1}+\mu_{2}+\mu_{3}$ \\
                            &$\chi_{2}^{B}$     &$\frac{1}{3}\left(2\mu_{1}+2\mu_{2}-\mu_{3}\right)$ \\
                            &$\chi_{3}^{B}$     &$\mu_{3}$ \\
                            &$C_{1}\chi_{2}^{B}+C_{2}\chi_{3}^{B}$
                            &$C_{1}^{2}\mu\left(\chi_{2}^{B}\right) + C_{2}^{2}\mu\left(\chi_{3}^{B}\right) + \frac{2}{\sqrt{3}}C_{1}C_{2}\left(\mu_{2}-\mu_{1}\right)$ \\
        $\left(q_{1}q_{2}\right)\left({\bar{q}}_{3}{\bar{q}}_{4}\right)$
                            &$\chi_{1}^{T}$     &$\mu_{1}+\mu_{2}+\mu_{3}+\mu_{4}$ \\
                            &$\chi_{2}^{T}$     &$\frac{1}{2}\left(\mu_{1}+\mu_{2}+\mu_{3}+\mu_{4}\right)$ \\
                            &$\chi_{3}^{T}$     &0 \\
                            &$\chi_{4}^{T}$     &$\mu_{1}+\mu_{2}$ \\
                            &$\chi_{5}^{T}$     &$\mu_{3}+\mu_{4}$ \\
                            &$\chi_{6}^{T}$     &0 \\
                            &$C_{1}\chi_{3}^{T}+C_{2}\chi_{6}^{T}$      &0 \\
                            &$C_{1}\chi_{5}^{T}+C_{2}\chi_{4}^{T}$
                            &$C_{1}^{2}\mu\left(\chi_{5}^{T}\right) + C_{2}^{2}\mu\left(\chi_{4}^{T}\right)$ \\
                            &$C_{1}\chi_{2}^{T}+C_{2}\chi_{4}^{T}+C_{3}\chi_{5}^{T}$
                            &$C_{1}^{2}\mu\left(\chi_{2}^{T}\right) + C_{2}^{2}\mu\left(\chi_{4}^{T}\right) + C_{3}^{2}\mu\left(\chi_{5}^{T}\right)
                            + \sqrt{2}C_{1}C_{2}\left(\mu_{3}-\mu_{4}\right) + \sqrt{2}C_{1}C_{3}\left(\mu_{2}-\mu_{1}\right)$ \\
                            &$C_{1}\chi_{2}^{T}+C_{2}\chi_{5}^{T}+C_{3}\chi_{4}^{T}$
                            &$C_{1}^{2}\mu\left(\chi_{2}^{T}\right) + C_{2}^{2}\mu\left(\chi_{5}^{T}\right) + C_{3}^{2}\mu\left(\chi_{4}^{T}\right)
                            + \sqrt{2}C_{1}C_{2}\left(\mu_{2}-\mu_{1}\right) + \sqrt{2}C_{1}C_{3}\left(\mu_{3}-\mu_{4}\right)$ \\
        \hline\hline
    \end{tabular}
\end{table*}


\subsection{Hadronic Properties}
Knowing the parameters $x_{i}$ and $R$ describing a hadronic state, one can subsequently calculate its charge radius and magnetic moment. Following the standard method, one can firstly evaluate the contribution of  a quark or an antiquark $i$ with electric charge $Q_{i}$ on charge radius \cite{DeGrand:1975cf}
\begin{equation}\label{rEi}
    \begin{aligned}
        {\left\langle r_{E}^{2} \right\rangle}_{i}&=
        Q_{i}R^{2} \frac{\alpha_{i} \left[2x_{i}^{2}\left(\alpha_{i}-1\right)+4\alpha_{i}+2\lambda_{i}-3\right]}
        {3x_{i}^{2} \left[2\alpha_{i}\left(\alpha_{i}-1\right)+\lambda_{i}\right]} \\
        &-Q_{i}R^{2} \frac{\lambda_{i} \left[4\alpha_{i}+2\lambda_{i}-2x_{i}^{2}-3\right]}
        {2x_{i}^{2} \left[2\alpha_{i}\left(\alpha_{i}-1\right)+\lambda_{i}\right]}
    \end{aligned}
\end{equation}
Then, one can directly sum up Eq.~(\ref{rEi}) to give the charge radius of a hadronic state \cite{Chodos:1974pn}
\begin{equation}\label{rEsum}
    r_{E}={\left| \sum\nolimits_{i} {\left\langle r_{E}^{2} \right\rangle}_{i} \right|}^{1/2}
\end{equation}
Note that Eq.~(\ref{rEsum}) also holds true for the chormomagneti-mixing systems having the identical quark constituents.

For evaluating magnetic moment, one can use the following equation \cite{DeGrand:1975cf,Wang:2016dzu}
\begin{equation}\label{mui}
    \mu_{i}=Q_{i}\bar{\mu}_{i}=Q_{i}\frac{R}{6} \frac{4\alpha_{i}+2\lambda_{i}-3}{2\alpha_{i}\left(\alpha_{i}-1\right)+\lambda_{i}}
\end{equation}
\begin{equation}\label{musum}
    \mu=\left\langle \psi_{spin} \left| \sum\nolimits_{i}g_{i}\mu_{i}S_{iz} \right| \psi_{spin} \right\rangle
\end{equation}
where $S_{iz}$ is the third component of spin for an individual quark or antiquark, and $g_{i}=2$ the Lander factor. If the chormomagneti-mixing enters so that the total spin wave function becomes
\begin{equation}
    \left| \psi_{spin} \right\rangle=C_{1}\chi_{1}+C_{2}\chi_{2}
\end{equation}
the magnetic moment becomes, corresondingly, 
\begin{equation}\label{mixingmu}
    \mu=C_{1}^{2}\mu\left(\chi_{1}\right)+C_{2}^{2}\mu\left(\chi_{2}\right)+2C_{1}C_{2}\mu^{\rm tr}\left(\chi_{1}, \chi_{2}\right)
\end{equation}
with $\mu^{\rm tr}$ the transition moment \cite{Bernotas:2012nz} and the set of $C_{1}$, $C_{2}$ the eigenvector for the mixing states.
Given the spin wave functions in Appendix B, we derive the magnetic moments for all the spin states and the mixing systems involved in this work, and list them collectively in Table \ref{tab:musum}.

Note that the transition moments in mixing systems are not symmetric under the exchange of the quarks $q_{1}$ and $q_{2}$ or $q_{3}$ and $q_{4}$ in the flavor wavefunctions, respectively. While isospin shall not be considered in constructing mass formula, it actually distinguishes the magnetic moments, between $\left(ud\right)s$ and $\left(du\right)s$ states for instance,
by the following wavefunction with isospin involved,  
\begin{equation}
    {\left| 1,0\right\rangle}_{I}=\frac{1}{\sqrt{2}} \left(ud+du\right) \quad
    {\left| 0,0\right\rangle}_{I}=\frac{1}{\sqrt{2}} \left(ud-du\right)
\end{equation}

Physically, one has to employ the above two states to evaluate magnetic moment, but the transition moments vanish eventually. Therefore, there is no transition moment between two mixing states, leaving only the terms with $C_{i}^{2}$ in Eq.~(\ref{mixingmu}).

Before ending of this section, we emphasize that there is a unique set of $x_{i}$ and $R$ for a given hadonic state with known flavor-spin-color structure and with certain CMI, and for those states carrying the same quantum number $J^{P}$, compared to others(Explain??). This eventually hinders us to find an simple formally-analytic mass formula in the case of the chromomagnetic-mixing systems since one has to perform variation analysis after diagonalizing CMI matrices. Keep this in mind while one one is calculating numerically the mass spectra and hadronic properties of mixing states, especially in the case of tetraquarks.


\section{Determination of Parameters}\label{Parameters}
In MIT bag model, the parameters quark mass $m_{i}$, bag constant $B$, coefficient $Z_{0}$ of ZPE and strong coupling constant $\alpha_{s}$ are determined based on the mass spectrum of the light hadrons $N$, $\Delta$, $\omega$ and $\Omega$ in their ground states. The results reads \cite{DeGrand:1975cf}
\begin{equation}\label{originalparas}
    \begin{Bmatrix}
        m_{n}=0\,{\rm GeV} & m_{s}=0.279\,{\rm GeV} \\
        Z_{0}=1.84 & B^{1/4}=0.145\,{\rm GeV} & \alpha_{s}=0.55
    \end{Bmatrix}
\end{equation}

We assume the values of the parameters in Eq.~(\ref{originalparas}) to be reasonable for both of light hadrons and heavy hadrons, with only exception that the strong coupling $\alpha_{s}$ can change weakly with the size $R$ of hadron around $0.55$, as given by Eq. (\ref{alphaS-mine}). In fact, one can check easily that the fixed value, near $0.55$ for instance, of $\alpha_{s}$ can lead to the mismatch error with respective to the measured masses of heavy mesons as large as $100$ MeV.

\renewcommand{\tabcolsep}{0.39cm}
\renewcommand{\arraystretch}{1.3}
\begin{table}[!htbp]
    \caption{Comparations of bag radius $R$ (${\rm GeV}^{-1}$) and bag mass $M_{bag}$ (${\rm GeV}$) between this work and original result for light hadrons.
    The experimental $M_{exp}$ (${\rm GeV}$) takes average of isospin.
    In this work, we choose $Z_{0}=1.83$ with little change to reduce errors of nucleon.}
    \label{tab:lighthadron}
    \begin{tabular}{cccccc}
        \hline\hline
        \multirow{2}{*}{State} & \multicolumn{2}{c}{\cite{Johnson:1975zp}} & \multicolumn{2}{c}{\rm this\ work} & \multirow{2}{*}{$M_{exp}$} \\
        \cline{2-3} \cline{4-5}
        & $R$ & $M_{bag}$ & $R$ & $M_{bag}$ \\ \hline
        $N$               &5.00      &0.938      &5.22      &0.932      &0.939  \\
        $\Delta$          &5.48      &1.233      &5.33      &1.241      &1.232  \\
        $\Lambda$         &4.95      &1.105      &5.26      &1.096      &1.116  \\
        $\Sigma$          &4.95      &1.144      &5.22      &1.137      &1.193  \\
        $\Xi$             &4.91      &1.289      &5.27      &1.282      &1.318  \\
        $\Sigma^{\ast}$   &5.43      &1.382      &5.38      &1.383      &1.385  \\
        $\Xi^{\ast}$      &5.39      &1.529      &5.42      &1.529      &1.533  \\
        $\Omega$          &5.35      &1.672      &5.46      &1.677      &1.672  \\
        $\pi$             &3.34      &0.280      &4.31      &0.348      &0.137  \\
        $\omega$          &4.71      &0.783      &4.55      &0.776      &0.783  \\
        $K$               &3.26      &0.497      &4.34      &0.561      &0.496  \\
        $K^{\ast}$        &4.65      &0.928      &4.63      &0.918      &0.894  \\
        $\phi$            &4.61      &1.068      &4.70      &1.064      &1.019  \\
        \hline\hline
    \end{tabular}
\end{table}

One remaining task is to determine the effective masses $m_{c}$ and $m_{b}$ of the charm and bottom quarks, as well as five kinds of binding energies $B_{QQ\prime}$. In this work, we choose the heavy-light mesons $D/D^{\ast}$ and $B/B^{\ast}$ to extract the heavy quark masses $m_{c}$ and $m_{b}$. Usage of the mesons $D^{\ast}$ and $B^{\ast}$ gives rise to
\begin{equation}\label{mcmb}
    \begin{Bmatrix}
        m_{c}=1.641\,{\rm GeV}, & m_{b}=5.093\,{\rm GeV}
    \end{Bmatrix}
\end{equation}
Substracting the heavy quark mass $m_{Q}$ from the measured meson mass spin-averaged, one can extract the residual masses to find the binding energy. Use of the mesons $D_{s}^{\ast}$, $B_{s}^{\ast}$, $J/\psi$, $\Upsilon$ and $B_{c}^{\ast}$ gives 
\begin{equation}\label{bingdings}
    \begin{Bmatrix}
        B_{cs}=-0.025\,{\rm GeV} & B_{cc}=-0.077\,{\rm GeV} \\
        B_{bs}=-0.032\,{\rm GeV} & B_{bb}=-0.128\,{\rm GeV} \\
        B_{bc}=-0.101\,{\rm GeV}
    \end{Bmatrix}
\end{equation}
for the color configuration $\boldsymbol{\bar{3}}$, Here in computation, we have used the mass $6.332$ GeV of $B_{c}^{\ast}$ predicted in Ref.~\cite{Ebert:2002pp}, due to the absence of measured mass for $B_{c}^{\ast}$.

In this work, the binding energy $B_{QQ^{\prime}}$ is determined for color $\boldsymbol{\bar{3}}$, which scaled, by Eq.~(\ref{colorfc}) through color configurations(???). For color $\boldsymbol{1}$, we have $2B_{QQ^{\prime}}$, and have the binding $-B_{QQ^{\prime}}/2$ for the color $\boldsymbol{6}$. In addition, the binding between quark and antiquark within tetraquark are $5B_{QQ^{\prime}}/4$ and $B_{QQ^{\prime}}/2$ for $\boldsymbol{6}\otimes\boldsymbol{\bar{6}}$ and $\boldsymbol{\bar{3}}\otimes\boldsymbol{3}$ respectively. Their scaled values can be found in Eq.~(\ref{cfcT}) (???).

In Tables \ref{tab:lightmeson}, \ref{tab:heavymeson}, \ref{tab:lightbaryon},  \ref{tab:1heavybaryon} and Table \ref{tab:23heavybaryon}, we list the computed masses, magnetic moments and charge radii for all established hadrons in their lowest-lying states without mixing. The results for the doubly and triply heavy baryons are also presented. Some remarks are in order: 

1. Only the well-established hadrons are used in this work and are enough to fix the model parameters. Some of the computed masses $M_{bag}$ deviates from the measured masses at the level of errors about $30\sim 40$ GeV, within the error range of MIT bag model, as shown in Table \ref{tab:lighthadron}.

2. The anti-particle mesons are not listed in Table \ref{tab:lightmeson} and \ref{tab:heavymeson} because they share the same mass, charge radius but minus magnetic moment in comparison with the mesens in Tables. For the same reason, we have ignored the anti-particle of the tetraquarks considered as well.

3. In Table \ref{tab:23heavybaryon}, our mass predictions $3.604$ GeV for the doubly charm baryon $\Xi_{cc}$ and $3.714$ GeV for $\Xi_{cc}^{\ast}$, are comparable to the quark model prediction $M\left(\Xi_{cc}^{\ast}\right)=3.727$ GeV ~\cite{Ebert:2004ck}, and also to (???) by the lattice QCD ~\cite{}.

4. In Table \ref{tab:lightbaryon}, the predicted magnetic moments are in good agreement with the measured values, from which we could further calculate the magnetic moments of heavy baryons and tetraquarks well.

5. Our prediction $0.75\,{\rm fm}$ for the proton charge radius is lower than the measured value $0.84\,{\rm fm}$, similar to original MIT bag model \cite{Johnson:1975zp,DeGrand:1975cf}. 

\renewcommand{\tabcolsep}{0.37cm}
\renewcommand{\arraystretch}{1.3}
\begin{table}[!htbp]
    \caption{Predicted spectra of ground-state light mesons. $\mu_{p}$ is theoretical magnetic moment of proton. Charge radius is in unit of fm.
    The blank cells follow above since MIT bag model confuses with isospin.}
    \label{tab:lightmeson}
    \begin{tabular}{cccccc}
        \hline\hline
        {\rm State}         &$R_{0}$    &$M_{bag}$  &$M_{exp}$  &$\mu_{bag}/\mu_{p}$    &$r_{E}$ \\ \hline
        $\pi^{+}$ 	        &4.31	    &0.348	    &0.137	    &-	                    &0.62 \\
        $\omega$ 	        &4.55	    &0.776	    &0.783	    &0                      &0 \\
        $K^{+}$ 	        &4.34	    &0.561	    &0.494	    &-                      &0.61 \\
        $K^{0}$ 		    &           &           &0.498	    &-	                    &0.13 \\
        $K^{\ast +}$ 	    &4.63	    &0.918	    &0.892	    &0.82                   &0.65 \\
        $K^{\ast 0}$ 	    &           &           &0.896	    &-0.06	                &0.15 \\
        $\phi$ 	            &4.70	    &1.064	    &1.019	    &0                      &0 \\
        \hline\hline
    \end{tabular}
\end{table}

\renewcommand{\tabcolsep}{0.3cm}
\renewcommand{\arraystretch}{1.3}
\begin{table}[!htbp]
    \caption{Predicted spectra of ground-state heavy mesons. The blank cell follows $5.325\,{\rm GeV}$ because it doesn’t identify specific state.}
    \label{tab:heavymeson}
    \begin{tabular}{cccccc}
        \hline\hline
        {\rm State}             &$R_{0}$    &$M_{bag}$      &$M_{exp}$  &$\mu_{bag}/\mu_{p}$    &$r_{E}$ \\ \hline
        $D^{+}$ 	            &3.63	    &1.835	        &1.870	    &-                      &0.46 \\
        $D^{0}$ 			    &           &               &1.865	    &-                      &0.25 \\
        $D^{\ast +}$ 	        &4.09	    &2.009(fitted)  &2.010	    &0.43                   &0.51 \\
        $D^{\ast 0}$ 		    &           &               &2.007	    &-0.35                  &0.29 \\
        $B^{+}$ 	            &3.14	    &5.248	        &5.279	    &-                      &0.42 \\
        $B^{0}$ 			    &           &               &5.280	    &-                      &0.17 \\
        $B^{\ast +}$ 	        &3.47	    &5.325(fitted)  &5.325      &0.47                   &0.46 \\
        $B^{\ast 0}$ 		    &           &               &           &-0.19                  &0.19 \\
        $D_{s}^{+}$ 	        &3.77	    &1.961	        &1.968	    &-                      &0.46 \\
        $D_{s}^{\ast +}$ 	    &4.17	    &2.112(fitted)  &2.112	    &0.39                   &0.51 \\
        $B_{s}^{0}$ 	        &3.35	    &5.346	        &5.367	    &-                      &0.16 \\
        $B_{s}^{\ast 0}$ 	    &3.62	    &5.415(fitted)  &5.415	    &0.36                   &0.17 \\
        $\eta_{c}$ 	            &3.15	    &3.002	        &2.984	    &-                      &0 \\
        $J/\psi$ 	            &3.54	    &3.097(fitted)  &3.097	    &0                      &0 \\
        $B_{c}^{+}$ 	        &2.53	    &6.274	        &6.275	    &-                      &0.29 \\
        $B_{c}^{\ast +}$ 	    &2.81	    &6.332(fitted)  &6.332	    &0.19                   &0.32 \\
        $\eta_{b}$ 	            &1.59	    &9.397	        &9.399	    &-                      &0 \\
        $\Upsilon$ 	            &1.80	    &9.460(fitted)  &9.460	    &0                      &0 \\
        \hline\hline
    \end{tabular}
\end{table}

\renewcommand{\tabcolsep}{0.25cm}
\renewcommand{\arraystretch}{1.3}
\begin{table}[!htbp]
    \caption{Predicted spectra of ground-state light baryons. $\mu_{p}$ dividing $\mu_{bag}$ is theoretical value and another one beside is experimental data.}
    \label{tab:lightbaryon}
    \begin{tabular}{ccccccc}
        \hline\hline
        {\rm State}         &$R_{0}$    &$M_{bag}$  &$M_{exp}$  &$\mu_{bag}/\mu_{p}$    &$\mu_{exp}/\mu_{p}$    &$r_{E}$ \\ \hline
        $p$ 	            &5.22	    &0.932	    &0.938	    &1	                    &1	                    &0.75 \\
        $n$ 			    &           &           &0.940	    &-0.67	                &-0.69	                &0 \\
        $\Delta^{++}$ 	    &5.33	    &1.241	    &1.231	    &2.04	                &2.20	                &1.08 \\
        $\Delta^{+}$ 	    &           &           &1.235	    &1.02	                &0.97	                &0.77 \\
        $\Delta^{0}$ 		&           &           &1.233	    &0	                    &-	                    &0 \\
        $\Lambda$ 	        &5.26	    &1.096	    &1.116	    &-0.26	                &-0.22	                &0.18 \\
        $\Sigma^{+}$ 	    &5.22	    &1.137	    &1.189	    &0.98	                &0.88	                &0.77 \\
        $\Sigma^{0}$ 		&           &           &1.193	    &0.31	                &-	                    &0.17 \\
        $\Sigma^{-}$ 		&           &           &1.197	    &-0.36	                &-0.42	                &0.73 \\
        $\Sigma^{\ast +}$ 	&5.38	    &1.383	    &1.383	    &1.11	                &-	                    &0.79 \\
        $\Sigma^{\ast 0}$   &           &           &1.384	    &0.08	                &-	                    &0.18 \\
        $\Sigma^{\ast -}$ 	&           &           &1.387	    &-0.95	                &-	                    &0.75 \\
        $\Xi^{0}$ 	        &5.27       &1.282	    &1.315	    &-0.57	                &-0.45	                &0.25 \\
        $\Xi^{-}$ 			&           &           &1.322	    &-0.23	                &-0.23	                &0.72 \\
        $\Xi^{\ast 0}$ 	    &5.42	    &1.529	    &1.532	    &0.17	                &-	                    &0.26 \\
        $\Xi^{\ast -}$ 		&           &           &1.535	    &-0.87	                &-	                    &0.74 \\
        $\Omega$ 	        &5.46	    &1.677	    &1.672	    &-0.79	                &-0.72	                &0.72 \\
        \hline\hline
    \end{tabular}
\end{table}

\renewcommand{\tabcolsep}{0.37cm}
\renewcommand{\arraystretch}{1.3}
\begin{table}[!htbp]
    \caption{Predicted spectra of ground-state and non-mixing singly heavy baryons.}
    \label{tab:1heavybaryon}
    \begin{tabular}{cccccc}
        \hline\hline
        {\rm State}             &$R_{0}$    &$M_{bag}$  &$M_{exp}$  &$\mu_{bag}/\mu_{p}$    &$r_{E}$ \\ \hline
        $\Lambda_{c}$ 	        &4.86	    &2.270	    &2.286	    &0.18	                &0.60 \\
        $\Sigma_{c}^{++}$ 	    &4.82	    &2.411	    &2.454	    &0.76	                &0.92 \\
        $\Sigma_{c}^{+}$ 		&           &           &2.453	    &0.15	                &0.60 \\
        $\Sigma_{c}^{0}$ 		&           &           &2.454	    &-0.47	                &0.35 \\
        $\Sigma_{c}^{\ast ++}$ 	&5.01	    &2.512	    &2.518	    &1.46	                &0.95 \\
        $\Sigma_{c}^{\ast +}$ 	&           &           &2.518	    &0.50	                &0.62 \\
        $\Sigma_{c}^{\ast 0}$ 	&           &           &2.518	    &-0.46	                &0.37 \\
        $\Omega_{c}$ 	        &4.93	    &2.680	    &2.695	    &-0.38                  &0.28 \\
        $\Omega_{c}^{\ast}$ 	&5.10	    &2.764	    &2.766	    &-0.32                  &0.29 \\
        $\Lambda_{b}$ 	        &4.60	    &5.648	    &5.620	    &-0.03	                &0.25 \\
        $\Sigma_{b}^{+}$ 	    &4.64	    &5.835	    &5.811	    &0.80	                &0.72 \\
        $\Sigma_{b}^{0}$ 		&           &           &-	        &0.21	                &0.26 \\
        $\Sigma_{b}^{-}$ 	    &           &           &5.816	    &-0.39	                &0.62 \\
        $\Sigma_{b}^{\ast +}$ 	&4.73	    &5.872	    &5.830	    &1.18	                &0.73 \\
        $\Sigma_{b}^{\ast 0}$ 	&           &           &-	        &0.27	                &0.26 \\
        $\Sigma_{b}^{\ast -}$ 	&           &           &5.835	    &-0.64	                &0.63 \\
        $\Omega_{b}$ 	        &4.77	    &6.080	    &6.046	    &-0.31                  &0.60 \\
        $\Omega_{b}^{\ast}$ 	&4.84	    &6.112	    &-	        &-0.51                  &0.60 \\
        \hline\hline
    \end{tabular}
\end{table}

\renewcommand{\tabcolsep}{0.51cm}
\renewcommand{\arraystretch}{1.3}
\begin{table}[!htbp]
    \caption{Predicted spectra of ground-state and non-mixing doubly and triply heavy baryons.}
    \label{tab:23heavybaryon}
    \begin{tabular}{ccccc}
        \hline\hline
        {\rm State}             &$R_{0}$    &$M_{bag}$  &$\mu_{bag}/\mu_{p}$    &$r_{E}$ \\ \hline
        $\Xi_{cc}^{++}$ 	    &4.42	    &3.604	    &0.05                   &0.78 \\
        $\Xi_{cc}^{+}$ 			&           &           &0.33	                &0.45 \\
        $\Xi_{cc}^{\ast ++}$ 	&4.64	    &3.714	    &0.94	                &0.82 \\
        $\Xi_{cc}^{\ast +}$ 	&           &           &0.06	                &0.47 \\
        $\Xi_{bb}^{0}$ 	        &3.71	    &10.311	    &-0.20	                &0.29 \\
        $\Xi_{bb}^{-}$ 			&           &           &0.04	                &0.45 \\
        $\Xi_{bb}^{\ast 0}$ 	&3.87	    &10.360	    &0.43	                &0.30 \\
        $\Xi_{bb}^{\ast -}$ 	&           &           &-0.31	                &0.47 \\
        $\Omega_{cc}$ 	        &4.49	    &3.726	    &0.31	                &0.48 \\
        $\Omega_{cc}^{\ast}$ 	&4.69	    &3.820	    &0.12	                &0.50 \\
        $\Omega_{bb}$ 	        &3.83	    &10.408	    &0.03	                &0.45 \\
        $\Omega_{bb}^{\ast}$ 	&3.97	    &10.451	    &-0.27	                &0.47 \\
        $\Omega_{ccc}$ 	        &4.22	    &4.841	    &0.52	                &0.68 \\
        $\Omega_{ccb}$ 	        &3.75	    &8.112	    &0.24	                &0.44 \\
        $\Omega_{ccb}^{\ast}$ 	&3.83	    &8.133	    &0.31	                &0.45 \\
        $\Omega_{cbb}$ 	        &3.18	    &11.373	    &-0.10	                &0.11 \\
        $\Omega_{cbb}^{\ast}$ 	&3.31	    &11.402	    &0.10	                &0.11 \\
        $\Omega_{bbb}$ 	        &2.59	    &14.626	    &-0.09	                &0.28 \\
        \hline\hline
    \end{tabular}
\end{table}


\section{Mixing States in baryons and Tetraquarks}\label{Hadrons}
\subsection{Baryons of Mixing Systems}
\renewcommand{\tabcolsep}{0.73cm}
\renewcommand{\arraystretch}{1.3}
\begin{table*}[!htbp]
    \caption{Predicted spectra of ground-state and mixing heavy baryons. The experimental data $M_{exp}$ takes average of isospin.
    Magnetic moments and charge radii are organized in the order of $I_{3}=\frac{1}{2},\,-\frac{1}{2}$ for $I=\frac{1}{2}$.}
    \label{tab:mixingbaryon}
    \begin{tabular}{ccccccc}
        \hline\hline
        {\rm State}         &{\rm Eigenvector}  &$R_{0}$    &$M_{bag}$  &$M_{exp}$  &$\mu_{bag}/\mu_{p}$    &$r_{E}$ \\ \hline
        \multirow{2}{*}{$\left(\Xi_{c}^{\prime},\ \Xi_{c}\right)$}
                            &(0.05, 1.00)       &4.89	    &2.436	    &2.470      &0.18, 0.18             &0.63, 0.32 \\
                            &(-1.00, 0.05)      &4.88	    &2.544	    &2.579      &0.20, -0.43            &0.63, 0.32 \\
        $\Xi_{c}^{\ast}$ 	&1.00               &5.06	    &2.636	    &2.646      &0.58, -0.39            &0.65, 0.33 \\
        \multirow{2}{*}{$\left(\Xi_{b}^{\prime},\ \Xi_{b}\right)$}
                            &(0.01, 1.00)       &4.64	    &5.805	    &5.795      &-0.03, -0.03           &0.30, 0.60 \\
                            &(-1.00, 0.01)      &4.71	    &5.956	    &5.935      &0.25, -0.35            &0.30, 0.61 \\
        $\Xi_{b}^{\ast}$ 	&1.00               &4.79	    &5.991	    &5.954      &0.34, -0.58            &0.31, 0.62 \\
        \multirow{2}{*}{$\left(\Xi_{bc}^{\prime},\ \Xi_{bc}\right)$}
                            &(0.39, 0.92)       &4.22	    &7.015	    &-          &0.45, -0.20            &0.58, 0.19 \\
                            &(-0.92, 0.39)      &4.09	    &6.953	    &-          &0.01, 0.12             &0.56, 0.19 \\
        $\Xi_{bc}^{\ast}$ 	&1.000              &4.31	    &7.044	    &-          &0.69, -0.13            &0.59, 0.20 \\
        \multirow{2}{*}{$\left(\Omega_{bc}^{\prime},\ \Omega_{bc}\right)$}
                            &(0.40, 0.92)       &4.29	    &7.116	    &-          &-0.16                  &0.15 \\
                            &(-0.92, 0.40)      &4.18	    &7.064	    &-          &0.11                   &0.14 \\
        $\Omega_{bc}^{\ast}$&1.000              &4.37	    &7.142	    &-          &-0.08                  &0.15 \\
        \hline\hline
    \end{tabular}
\end{table*}

\renewcommand{\tabcolsep}{0.62cm}
\renewcommand{\arraystretch}{1.3}
\begin{table*}[!htbp]
    \caption{Predicted spectra of singly heavy tetraquarks $nn\bar{s}\bar{c}$ and $nn\bar{s}\bar{b}$.
    Magnetic moments and charge radii are organized in the order of $I_{3}=1,\,0,\,-1$ for $I=1$.}
    \label{tab:1heavytetraquark}
    \begin{tabular}{ccccccc}
        \hline\hline
        {\rm State}     &$J^{P}$    &{\rm Eigenvector}      &$R_{0}$    &$M_{bag}$  &$\mu_{bag}/\mu_{p}$    &$r_{E}$ \\ \hline
        \multirow{2}{*}{${(nn\bar{s}\bar{c})}^{I=1}$}
                        &\multirow{2}{*}{$0^{+}$}
                                    &(0.54, 0.84)           &5.73	    &3.218	    &-                      &0.91, 0.38, 0.73 \\
                        &           &(-0.84, 0.55)          &5.39	    &2.776	    &-                      &0.85, 0.35, 0.69 \\
                        &\multirow{3}{*}{$1^{+}$}
                                    &(0.81, 0.58, 0.10)     &5.46	    &3.001	    &0.95, 0.26, -0.43      &0.86, 0.36, 0.70 \\
                        &           &(0.25, -0.49, 0.84)    &5.57	    &3.154	    &0.44, 0.16, -0.13      &0.88, 0.37, 0.71 \\
                        &           &(-0.54, 0.65, 0.54)    &5.38	    &2.846	    &0.60, 0.43, -0.57      &0.85, 0.35, 0.69 \\
                        &$2^{+}$    &1.00                   &5.64	    &3.075	    &1.53, 0.45, -0.63      &0.89, 0.37, 0.72 \\
        \multirow{2}{*}{${(nn\bar{s}\bar{c})}^{I=0}$}
                        &\multirow{2}{*}{$0^{+}$}
                                    &(0.63, 0.77)           &5.56	    &2.934	    &-                      &0.37 \\
                        &           &(-0.78, 0.63)          &5.19	    &2.513	    &-                      &0.34 \\
                        &\multirow{3}{*}{$1^{+}$}
                                    &(0.77, 0.07, 0.64)     &5.40	    &2.895	    &0.16                   &0.35 \\
                        &           &(-0.23, -0.90, 0.36)   &5.57	    &3.056	    &0.31                   &0.37 \\
                        &           &(-0.60, 0.43, 0.68)    &5.35	    &2.674	    &0.17                   &0.35 \\
                        &$2^{+}$    &1.00                   &5.66	    &3.063	    &0.45                   &0.37 \\

        \multirow{2}{*}{${(nn\bar{s}\bar{b})}^{I=1}$}
                        &\multirow{2}{*}{$0^{+}$}
                                    &(0.53, 0.85)           &5.53	    &6.580	    &-                      &1.07, 0.71, 0.36 \\
                        &           &(-0.85, 0.53)          &5.28	    &6.202	    &-                      &1.02, 0.68, 0.34 \\
                        &\multirow{3}{*}{$1^{+}$}
                                    &(0.66, 0.75, 0.06)     &5.35	    &6.419	    &1.13, 0.33, -0.47      &1.03, 0.69, 0.35 \\
                        &           &(0.36, -0.38, 0.85)    &5.43	    &6.554	    &0.52, 0.30, 0.09       &1.05, 0.70, 0.35 \\
                        &           &(-0.66, 0.55, 0.52)    &5.22	    &6.228	    &0.83, 0.31, -0.21      &1.01, 0.67, 0.34 \\
                        &$2^{+}$    &1.00                   &5.47	    &6.446	    &1.69, 0.64, -0.40      &1.05, 0.70, 0.36 \\
        \multirow{2}{*}{${(nn\bar{s}\bar{b})}^{I=0}$}
                        &\multirow{2}{*}{$0^{+}$}
                                    &(0.66, 0.75)           &5.41	    &6.327	    &-                      &0.69 \\
                        &           &(-0.75, 0.66)          &5.14	    &5.980	    &-                      &0.66 \\
                        &\multirow{3}{*}{$1^{+}$}
                                    &(0.66, -0.23, 0.72)    &5.33	    &6.322	    &0.30                   &0.68 \\
                        &           &(-0.46, -0.88, 0.14)   &5.38	    &6.454	    &0.34                   &0.69 \\
                        &           &(-0.60, 0.43, 0.68)    &5.14	    &6.038	    &0.30                   &0.66 \\
                        &$2^{+}$    &1.00                   &5.48	    &6.431	    &0.64                   &0.70 \\
        \hline\hline
    \end{tabular}
\end{table*}

\renewcommand{\tabcolsep}{0.59cm}
\renewcommand{\arraystretch}{1.2}
\begin{table*}[!htbp]
    \caption{Predicted spectra of doubly heavy tetraquarks $nn\bar{c}\bar{c}$, $nn\bar{b}\bar{b}$ and $nn\bar{c}\bar{b}$.
    Magnetic moments and charge radii are organized in the order of $I_{3}=1,\,0,\,-1$ for $I=1$.}
    \label{tab:2heavytetraquark1}
    \begin{tabular}{ccccccc}
        \hline\hline
        {\rm State}     &$J^{P}$    &{\rm Eigenvector}      &$R_{0}$    &$M_{bag}$  &$\mu_{bag}/\mu_{p}$    &$r_{E}$ \\ \hline
        \multirow{2}{*}{${(nn\bar{c}\bar{c})}^{I=1}$}
                        &\multirow{2}{*}{$0^{+}$}
                                    &(0.40, 0.92)           &5.40	    &4.342	    &-                      &0.56, 0.54, 0.94 \\
                        &           &(-0.91, 0.41)          &5.04	    &4.032	    &-                      &0.52, 0.50, 0.88 \\
                        &$1^{+}$    &1.00                   &5.22	    &4.117	    &0.49, -0.02, -0.51     &0.54, 0.52, 0.91 \\
                        &$2^{+}$    &1.00                   &5.32	    &4.179	    &1.00, -0.02, -1.04     &0.55, 0.53, 0.93 \\
        \multirow{2}{*}{${(nn\bar{c}\bar{c})}^{I=0}$}
                        &\multirow{2}{*}{$1^{+}$}
                                    &(0.97, 0.25)           &5.15	    &3.925	    &-0.32                  &0.51 \\
                        &           &(-0.24, 0.97)          &5.30	    &4.205	    &0.30                   &0.53 \\
        \multirow{2}{*}{${(nn\bar{b}\bar{b})}^{I=1}$}
                        &\multirow{2}{*}{$0^{+}$}
                                    &(0.17, 0.99)           &4.90	    &11.093	    &-                      &0.92, 0.59, 0.38 \\
                        &           &(-0.98, 0.18)          &4.77	    &10.834	    &-                      &0.90, 0.58, 0.37 \\
                        &$1^{+}$    &1.00                   &4.83	    &10.854	    &0.65, 0.18, -0.28      &0.91, 0.58, 0.38 \\
                        &$2^{+}$    &1.00                   &4.88	    &10.878	    &1.31, 0.37, -0.56      &0.92, 0.59, 0.38 \\
        \multirow{2}{*}{${(nn\bar{b}\bar{b})}^{I=0}$}
                        &\multirow{2}{*}{$1^{+}$}
                                    &(1.00, 0.08)           &4.76	    &10.654	    &0.06                   &0.57 \\
                        &           &(-0.08, 1.00)          &4.83	    &10.982	    &0.31                   &0.58 \\

        \multirow{2}{*}{${(nn\bar{c}\bar{b})}^{I=1}$}
                        &\multirow{2}{*}{$0^{+}$}
                                    &(0.30, 0.95)           &5.16	    &7.714	    &-                      &0.78, 0.25, 0.70 \\
                        &           &(-0.95, 0.31)          &4.91	    &7.438	    &-                      &0.74, 0.23, 0.67 \\
                        &\multirow{3}{*}{$1^{+}$}
                                    &(0.65, 0.76, 0.09)     &5.04	    &7.509	    &0.98, 0.22, -0.54      &0.76, 0.24, 0.68 \\
                        &           &(0.13, -0.22, 0.97)    &5.10	    &7.699	    &-0.07, -0.12, -0.18    &0.77, 0.24, 0.69 \\
                        &           &(-0.75, 0.61, 0.24)    &4.96	    &7.465	    &0.79, 0.16, -0.47      &0.75, 0.24, 0.67 \\
                        &$2^{+}$    &1.00                   &5.12	    &7.531	    &1.16, 0.18, -0.80      &0.78, 0.25, 0.69 \\
        \multirow{2}{*}{${(nn\bar{c}\bar{b})}^{I=0}$}
                        &\multirow{2}{*}{$0^{+}$}
                                    &(0.93, 0.37)           &4.96	    &7.502	    &-                      &0.24 \\
                        &           &(-0.38, 0.93)          &4.84	    &7.260	    &-                      &0.23 \\
                        &\multirow{3}{*}{$1^{+}$}
                                    &(0.91, -0.36, 0.21)    &4.93	    &7.518	    &0.10                   &0.23 \\
                        &           &(-0.39, -0.92, 0.10)   &5.07	    &7.605	    &0.28                   &0.24 \\
                        &           &(-0.16, 0.18, 0.97)    &4.90	    &7.288	    &-0.13                  &0.23 \\
                        &$2^{+}$    &1.00                   &5.14	    &7.483	    &0.18                   &0.25 \\
        \hline\hline
    \end{tabular}
\end{table*}

\renewcommand{\tabcolsep}{0.84cm}
\renewcommand{\arraystretch}{1.2}
\begin{table*}[!htbp]
    \caption{Predicted spectra of doubly heavy tetraquarks $ss\bar{c}\bar{c}$, $ss\bar{b}\bar{b}$ and $ss\bar{c}\bar{b}$.}
    \label{tab:2heavytetraquark0}
    \begin{tabular}{ccccccc}
        \hline\hline
        {\rm State}     &$J^{P}$    &{\rm Eigenvector}      &$R_{0}$    &$M_{bag}$  &$\mu_{bag}/\mu_{p}$    &$r_{E}$ \\ \hline
        \multirow{2}{*}{$ss\bar{c}\bar{c}$}
                        &\multirow{2}{*}{$0^{+}$}
                                    &(0.49, 0.87)           &5.48	    &4.521	    &-                      &0.92 \\
                        &           &(-0.87, 0.50)          &5.13	    &4.299	    &-                      &0.87 \\
                        &$1^{+}$    &1.00                   &5.30	    &4.382	    &-0.44                  &0.89 \\
                        &$2^{+}$    &1.00                   &5.39	    &4.433	    &-0.88                  &0.91 \\
        \multirow{2}{*}{$ss\bar{b}\bar{b}$}
                        &\multirow{2}{*}{$0^{+}$}
                                    &(0.24, 0.97)           &5.01	    &11.232	    &-                      &0.32 \\
                        &           &(-0.97, 0.25)          &4.88	    &11.078	    &-                      &0.31 \\
                        &$1^{+}$    &1.00                   &4.94	    &11.098	    &-0.21                  &0.31 \\
                        &$2^{+}$    &1.00                   &4.98	    &11.119	    &-0.43                  &0.32 \\
        \multirow{2}{*}{$ss\bar{c}\bar{b}$}
                        &\multirow{2}{*}{$0^{+}$}
                                    &(0.40, 0.92)           &5.26	    &7.875	    &-                      &0.67 \\
                        &           &(-0.91, 0.40)          &5.01	    &7.693	    &-                      &0.64 \\
                        &\multirow{3}{*}{$1^{+}$}
                                    &(0.70, 0.71, 0.11)     &4.98	    &7.757	    &-0.41                  &0.64 \\
                        &           &(0.17, -0.32, 0.93)    &5.00	    &7.858	    &-0.17                  &0.64 \\
                        &           &(-0.69, 0.63, 0.35)    &4.88	    &7.716	    &-0.24                  &0.62 \\
                        &$2^{+}$    &1.00                   &5.20	    &7.779	    &-0.66                  &0.66 \\
        \hline\hline
    \end{tabular}
\end{table*}

\renewcommand{\tabcolsep}{0.74cm}
\renewcommand{\arraystretch}{1.3}
\begin{table*}[!htbp]
    \caption{Predicted spectra of doubly heavy tetraquarks $ns\bar{c}\bar{c}$, $ns\bar{b}\bar{b}$.
    Magnetic moments and charge radii are organized in the order of $I_{3}=\frac{1}{2},\,-\frac{1}{2}$ for $I=\frac{1}{2}$.}
    \label{tab:2heavytetraquark1/2}
    \begin{tabular}{ccccccc}
        \hline\hline
        {\rm State}     &$J^{P}$    &{\rm Eigenvector}      &$R_{0}$    &$M_{bag}$  &$\mu_{bag}/\mu_{p}$    &$r_{E}$ \\ \hline
        \multirow{2}{*}{$ns\bar{c}\bar{c}$}
                        &\multirow{2}{*}{$0^{+}$}
                                    &(0.44, 0.90)           &5.44	    &4.429	    &-                      &0.51, 0.93 \\
                        &           &(-0.89, 0.45)          &5.09	    &4.165	    &-                      &0.48, 0.88 \\
                        &\multirow{3}{*}{$1^{+}$}
                                    &(0.99, -0.07, 0.09)    &5.16	    &4.247	    &0.03, -0.47            &0.49, 0.89 \\
                        &           &(-0.11, -0.28, 0.95)   &5.23	    &4.314	    &0.35, -0.57            &0.49, 0.90 \\
                        &           &(0.04, 0.96, 0.29)     &5.04	    &4.091	    &-0.29, -0.37           &0.48, 0.87 \\
                        &$2^{+}$    &1.00                   &5.36	    &4.305	    &0.07, -0.96            &0.50, 0.92 \\
        \multirow{2}{*}{$ns\bar{b}\bar{b}$}
                        &\multirow{2}{*}{$0^{+}$}
                                    &(0.20, 0.98)           &4.96	    &11.160	    &-                      &0.62, 0.35 \\
                        &           &(-0.98, 0.20)          &4.83	    &10.955	    &-                      &0.60, 0.34 \\
                        &\multirow{3}{*}{$1^{+}$}
                                    &(1.00, -0.01, 0.03)    &4.79	    &10.974	    &0.22, -0.24            &0.60, 0.34 \\
                        &           &(-0.03, -0.09, 1.00)   &4.77	    &11.068	    &0.37, -0.54            &0.60, 0.34 \\
                        &           &(0.01, 1.00, 0.10)     &4.66	    &10.811	    &0.06, 0.06             &0.58, 0.33 \\
                        &$2^{+}$    &1.00                   &4.93	    &10.997	    &0.45, -0.50            &0.62, 0.35 \\
        \hline\hline
    \end{tabular}
\end{table*}

\renewcommand{\tabcolsep}{0.25cm}
\renewcommand{\arraystretch}{1.3}
\begin{table*}[!htbp]
    \caption{Comparison of predicted spectra via different methods for double heavy tetraquarks. Refs.\cite{Luo:2017eub} and \cite{Lu:2020rog} both apply CMI model.}
    \label{tab:comparison}
    \begin{tabular}{ccccccc}
        \hline\hline
        {\rm State}     &$J$    &{\rm This\ work}       &\cite{Luo:2017eub} {\rm Scheme\ 1}  &\cite{Luo:2017eub} {\rm Scheme\ 2}  &\cite{Lu:2020rog}      &\cite{Ebert:2007rn} \\ \hline
        ${(nn\bar{c}\bar{c})}^{I=1}$
                        &0      &4.032/4.342            &4.078/4.356                         &3.850/4.128                         &4.195/4.414            &4.056 \\
                        &1      &4.117                  &4.201                               &3.973                               &4.268                  &4.079 \\
                        &2      &4.179                  &4.271                               &4.044                               &4.318                  &4.118 \\
        ${(nn\bar{c}\bar{c})}^{I=0}$
                        &1      &3.925/4.205            &4.007/4.204                         &3.779/3.977                         &4.041/4.313            &3.935 \\
        ${(nn\bar{b}\bar{b})}^{I=1}$
                        &0      &10.834/11.093          &10.841/10.937                       &10.637/10.734                       &10.765/11.019          &10.648 \\
                        &1      &10.854                 &10.875                              &10.671                              &10.779                 &10.657 \\
                        &2      &10.878                 &10.897                              &10.694                              &10.799                 &10.673 \\
        ${(nn\bar{b}\bar{b})}^{I=0}$
                        &1      &10.654/10.982          &10.686/10.821                       &10.483/10.617                       &10.550/10.951          &10.502 \\
        ${(nn\bar{c}\bar{b})}^{I=1}$
                        &0      &7.438/7.714            &7.457/7.643                         &7.241/7.428                         &7.519/7.740            &7.383 \\
                        &1      &7.465/7.509/7.699      &7.473/7.548/7.609                   &7.258/7.332/7.393                   &7.537/7.561/7.729      &7.396/7.403 \\
                        &2      &7.531                  &7.582                               &7.367                               &7.586                  &7.422 \\
        ${(nn\bar{c}\bar{b})}^{I=0}$
                        &0      &7.260/7.502            &7.256/7.429                         &7.041/7.213                         &7.297/7.580            &7.239 \\
                        &1      &7.288/7.518/7.605      &7.321/7.431/7.516                   &7.106/7.215/7.301                   &7.325/7.607/7.666      &7.246 \\
                        &2      &7.483                  &7.530                               &7.315                               &7.697                  &- \\
        $ns\bar{c}\bar{c}$
                        &0      &4.165/4.429            &4.236/4.514                         &3.933/4.210                         &4.323/4.512            &4.221 \\
                        &1      &4.091/4.247/4.314      &4.225/4.363/4.400                   &3.921/4.060/4.096                   &4.232/4.427/4.512      &4.143/4.239 \\
                        &2      &4.305                  &4.434                               &4.131                               &4.440                  &4.271 \\
        $ns\bar{b}\bar{b}$
                        &0      &10.955/11.160          &10.999/11.095                       &10.707/10.804                       &10.883/11.098          &10.802 \\
                        &1      &10.811/10.974/11.068   &10.911/11.010/11.037                &10.619/10.718/10.745                &10.734/10.897/11.046   &10.706/10.809 \\
                        &2      &10.997                 &11.060                              &10.769                              &10.915                 &10.823 \\
        $ss\bar{c}\bar{c}$
                        &0      &4.299/4.521            &4.395/4.672                         &4.016/4.293                         &4.417/4.587            &4.359 \\
                        &1      &4.382                  &4.526                               &4.146                               &4.493                  &4.375 \\
                        &2      &4.433                  &4.597                               &4.218                               &4.536                  &4.402 \\
        $ss\bar{b}\bar{b}$
                        &0      &11.078/11.232          &11.157/11.254                       &10.777/10.875                       &10.972/11.155          &10.932 \\
                        &1      &11.098                 &11.199                              &10.820                              &10.986                 &10.939 \\
                        &2      &11.119                 &11.224                              &10.844                              &11.004                 &10.950 \\
        $ss\bar{c}\bar{b}$
                        &0      &7.693/7.875            &7.774/7.960                         &7.394/7.581                         &7.735/7.894            &7.673 \\
                        &1      &7.716/7.757/7.858      &7.793/7.872/7.924                   &7.414/7.493/7.545                   &7.752/7.775/7.881      &7.683/7.684 \\
                        &2      &7.779                  &7.908                               &7.529                               &7.798                  &7.701 \\
        \hline\hline
    \end{tabular}
\end{table*}

Recall a chromomagnetic mixing, in which a diquark or antiquark with different flavors, beyond $SU(2)$ isospin symmetry, can be symmetric or asymmetric in flavor space. To begin with, we first consider the mixed systems $(\phi^{B}\chi_{2}^{B},\phi^{B}\chi_{3}^{B})$ for baryons with $J^{P}={\frac{1}{2}}^{+}$, for which the baryons $\Xi_{c}$ , $\Xi_{b}$, $\Xi_{bc}$ and $\Omega_{bc}$ are involved.

Given the matrices of CMI calculated with the wavefunctions in color and spin space(Appendix A and B), one can evaluate(Appendix C) the CMI matrices in the chromomagnetic-mixed system, where the subscripts of $C_{ij}$ denote the flavor configurations. We take a quark-diquark configurations $(ns)c$ for $\Xi_{c}$, $(ns)b$ for $\Xi_{b}$, $(bc)n$ for $\Xi_{bc}$ and $(bc)s$ for $\Omega_{bc}$. In addition, we add the binding energies $B_{cs}$, $B_{bs}$ and $B_{bc}$ to the shifted mass $\Delta M$ directly. 

First of all, one has to diagonalize the matrix (\ref{mixing23}) to get two eigenvalues to find a mass formula. Once $x_{i}$, $R$ are determined, the eigenvectors of the matrix (\ref{mixing23}) can be solved by variation numerically and can thereby be used to identify and denote the mixing states.

In Table \ref{tab:mixingbaryon}, we list our computed results of masses and other properties for the mixed system of heavy baryons considered. As shown in the second column, the net effects of the state-mixing are not so significant, even negligible in the case of singly heavy baryons. Thus, the masses of the mixed states are comparable with the measured masses of $\Xi_{c}$, $\Xi_{c}^{\prime}$, $\Xi_{b}$ and $\Xi_{b}^{\prime}$, as seen in the fifth column with reasonable errors.

On the other hand, our results indicate that the mixing effects becomes weaker as charm quark $c$ is replaced by bottom quark $b$ in $(ns)c$, or the strange quark $s$ replaced by $n$ quarks in $(bc)s$. This can be due to breaking of heavy quark symmetry and the $1/m_{Q}$ suppression of the off-diagonal elements of the matrix (\ref{mixing23}).


\subsection{Singly Heavy Tetraquarks}
Let us consider singly heavy tetraquarks $nn\bar{s}\bar{c}$ and $nn\bar{s}\bar{b}$. To find chromomagnetic mixing of them, we introduce the mixing systems such as
$\phi_{2}^{T}\chi_{3}^{T}-\phi_{1}^{T}\chi_{6}^{T}$ with $J^{P}=0^{+}$ and $I=1$,
$\phi_{1}^{T}\chi_{3}^{T}-\phi_{2}^{T}\chi_{6}^{T}$ with $J^{P}=0^{+}$ and $I=0$,
$\phi_{2}^{T}\chi_{2}^{T}-\phi_{2}^{T}\chi_{4}^{T}-\phi_{1}^{T}\chi_{5}^{T}$ with $J^{P}=1^{+}$ and $I=1$ and
$\phi_{1}^{T}\chi_{2}^{T}-\phi_{1}^{T}\chi_{4}^{T}-\phi_{2}^{T}\chi_{5}^{T}$ with $J^{P}=1^{+}$ and $I=0$.
Further, the binding energy matrices become diagonal in mass formula since the mixed states with two color configurations and spin states are orthogonal. Diagonalizing should be applied to the whole sum of the interaction matrices when calculating the hadron mass.

Following variational principle, we diagonalize the 2-by-2 matrix to solve two analytical eigenvalues and then construct the mass formula as usual.
The application of the same procedure on the 3-by-3 matrix is, however, not straightforward, for which the eigenvalues are some roots of a cubic equation.
For this, we choose to numerically scan three root sets of $x_{i}$ and $R$ to solve the cubic equation numerically so that one can obtain minimal masses with three root eigenvalues.

Our numerical results are shown in Table \ref{tab:1heavytetraquark}, with a notable tetraquark of a isosinglet $nn\bar{s}\bar{c}$ with $J^{P}=0^{+}$, having the masses $2.934$ GeV and $2.513$ GeV for two mixing states, respectively.
Considering that the measured mass was reported to be $2866\pm 7$ MeV \cite{LHCb:2020kd} and the quark model prediction is about $2863.4\pm 12$ MeV\cite{Karliner:2020vsi}, our mass prediction $2.934$ GeV is larger even if the model error $40$ MeV is subtracted. Instead of this, if we rather ignore the mixing effects and calculate the masses of the $\phi_{1}^{T}\chi_{3}^{T}$ and $\phi_{2}^{T}\chi_{6}^{T}$ states, as Karliner suggested for the color $\boldsymbol{\bar{3}}\otimes\boldsymbol{3}$ configuration, the resulted masses are around $2.7$ GeV and non-physical. 

On the other hand, the masses due to chromomagnetic mixing with $J^{P}=0^{+}$ leads to a large splitting, as well as the prediction \cite{Luo:2017eub}, where the same CMI model is adopted. This suggests that the open-heavy-flavor exotic $X_{0}(2900)$ is not likely to be isosinglet $nn\bar{s}\bar{c}$ tetraquark, unless a larger error $70$ MeV for the $K$ meson in Table \ref{tab:lightmeson} were accepted.


\subsection{Doubly Heavy Tetraquarks}
Finally, let us consider the doubly heavy tetraquarks with strange number $S\leq 2$. Possible mixing can occur in the systems such as $(\phi_{2}^{T}\chi_{3}^{T},\phi_{1}^{T}\chi_{6}^{T})$ with $J^{P}=0^{+}$ for isotriplet $nn\bar{Q}\bar{Q}$, $ns\bar{Q}\bar{Q}$ and $ss\bar{Q}\bar{Q}$,
$(\phi_{2}^{T}\chi_{5}^{T},\phi_{1}^{T}\chi_{4}^{T})$ with $J^{P}=1^{+}$ for isosinglet $nn\bar{Q}\bar{Q}$,
and $(\phi_{2}^{T}\chi_{2}^{T},\phi_{2}^{T}\chi_{5}^{T},\phi_{1}^{T}\chi_{4}^{T})$ with $J^{P}=1^{+}$ for $ns\bar{Q}\bar{Q}$. The case of the state
$nn\bar{c}\bar{b}$ follows the same case of $nn\bar{s}\bar{c}$. We neglect the mixing among six the spin-color states for the tetraquark $ns\bar{c}\bar{b}$, for which the sextic equations are required to solved directly.

We note that near-threshold effect is not accessible to address in this work owing to the uncertainty of hadron mass about $40$ MeV. Given none of confirmed tetraquark state available, we can only make predictions with error analysis. For the different conclusions about stability of isosinglet $nn\bar{b}\bar{b}$ with $J^{P}=1^{+}$, we refer to two calculations about chromomagnetic mixing for the doubly heavy tetraquarks \cite{Luo:2017eub,Lu:2020rog}. 

We use CMI model to identify the doubly heavy tetraquarks and compute the dynamical properties for these exotic hadrons. Our computed results for the doubly heavy tetraquarks are listed in Tables \ref{tab:2heavytetraquark1}, \ref{tab:2heavytetraquark0} and \ref{tab:2heavytetraquark1/2}. They are summarized in Table \ref{tab:comparison}, compared with other's works cited. In these works, Ref.~\cite{Luo:2017eub} has adopted the CMI model and Ref.~\cite{Lu:2020rog} used relativized quark model with CMI to compute the hadron spectra. 

Compared to the CMI model, the MIT bag model calculates $C_{ij}$ and then compute hadron masses, giving a different dynamical properties for hadronic states with the help of variational method. We also note that adding of binding energy among heavy quarks, stemming from chromoelectric interaction, being Coulomb-like interaction, reconciles the predictions for the light hadrons and the heavy sectors in this work. 

As stated in the previous sections, our results imply larger mass splittings among the chromomagnetic-mixed systems. For singly heavy tetraquarks, the mass splitting is about $400$ MeV. For the doubly heavy sectors, the splitting can be reduced to $300$ MeV. 

Comparing our predictions with Ref.~\cite{Lu:2020rog}, a notable phenomenon is that the mass distinctions between two calculations are of regular for charmed or bottom sectors. For instance, for the systems $nn\bar{c}\bar{c}$ and $ns\bar{c}\bar{c}$, our spectra are $150$ MeV smaller while they are $75$ MeV larger for the systems $nn\bar{b}\bar{b}$ and $ns\bar{b}\bar{b}$. In the systems of $ss\bar{c}\bar{c}$ , our mass predictions are $100$ MeV smaller but $100$ MeV larger for the systems $ss\bar{b}\bar{b}$ systems. These regular mass distinctions remind us of errors in established states, where bottom baryons are overestimated and charmed baryons are underestimated. Thus, it remains to make a systematic error analysis, staring from the light hadrons to heavy hadrons, and thereby giving reliable prediction for the heavy tetraquarks.


\section{Summary and Discussions}\label{Discussions}
In this work, we have estimated the mass spectra, magnetic moments and charge radii for all the ground-state mesons, baryons and some tetraquarks,
in MIT bag model with CMI model which describes intrinsical properties of hadronic states. The model parameters are found to be universal for light hadrons and heavy hadrons, most of which, including running strong coupling, are determined from light hadrons.

The calculations on mesons and baryons aim to prove it with a model error about $40$ MeV. Differing from previous works of bag model, the strong coupling is running from $0.4$ to $0.6$ when bag radius varies in $R=3\sim 6$ GeV$^{-1}$.

The binding energies, assumed to exist among strange and heavy quarks, can offset errors for heavy quarkonia and give close predictions of singly heavy baryons as $\Omega_{c}^{\ast}$ and doubly heavy sectors as $\Xi_{cc}$. This added binding energy can be explained by one-gluon-exchange between two relatively heavy quarks which locate at the center of hadron, and can be made to be universal for the whole hadron spectrum. 

In princple, the bag model can be applied to any hadronic states even glueball and quark star, with CMI matrices calculable explicitly. This helps to enrich our understanding on the structures and mass composition of hadron. We promote the error of $40$ MeV to tetraquarks and conclude that $X_{0}(2900)$ isn't an isosinglet $nn\bar{s}\bar{c}$ state whose distinction is up to $70$ MeV. 


\medskip
\textbf{ACKNOWLEDGMENTS}

W. Z thanks Xiang Liu, Si-Qiang Luo and Hong-Tao An for useful discussions.
D. J. is supported by the National Natural Science Foundation of China under the no. 11565023 and
the Feitian Distinguished Professor Program of Gansu (2014-2016).

\medskip
\textbf{NOTICE}

The Mathematica code used in the numerical evaluations is attached as a supplementary material.


\medskip
\section*{Appendix A}\label{AppendixA}
\setcounter{equation}{0}
\renewcommand{\theequation}{A\arabic{equation}}
For meson, baryon, and tetraquark systems labeled $(q_{1}\bar{q}_{2})$ $(q_{1}q_{2})q_{3}$ and $(q_{1}q_{2})(\bar{q}_{3}\bar{q}_{4})$ respectively,
the total color wave functions studied from SU(3) symmetry read
\begin{equation}\label{pcolorM}
    \phi^{M}=\frac{1}{\sqrt{3}} \left(r\bar{r}+g\bar{g}+b\bar{b} \right)
\end{equation}
\begin{equation}\label{pcolorB}
    \phi^{B}=\frac{1}{\sqrt{6}} \left(gbr-bgr+brg-rbg+rgb-grb \right)
\end{equation}
\begin{equation}\label{pcolorT}
    \begin{aligned}
        \phi_{1}^{T}&=\frac{1}{\sqrt{6}} \left(rr\bar{r}\bar{r}+gg\bar{g}\bar{g}+bb\bar{b}\bar{b} \right)
        +\frac{1}{2\sqrt{6}} \left(rb\bar{b}\bar{r}+br\bar{b}\bar{r} \right.\\
        &\left. +gr\bar{g}\bar{r}+rg\bar{g}\bar{r}+gb\bar{b}\bar{g}+bg\bar{b}\bar{g}+gr\bar{r}\bar{g}+rg\bar{r}\bar{g}+gb\bar{g}\bar{b} \right. \\
        &\left. +bg\bar{g}\bar{b}+rb\bar{r}\bar{b}+br\bar{r}\bar{b} \right) \\
        \phi_{2}^{T}&=\frac{1}{2\sqrt{3}} \left(rb\bar{b}\bar{r}-br\bar{b}\bar{r}-gr\bar{g}\bar{r}+rg\bar{g}\bar{r}+gb\bar{b}\bar{g}-bg\bar{b}\bar{g} \right.\\
        &\left. +gr\bar{r}\bar{g}-rg\bar{r}\bar{g}-gb\bar{g}\bar{b}+bg\bar{g}\bar{b}-rb\bar{r}\bar{b}+br\bar{r}\bar{b}\right)
    \end{aligned}
\end{equation}
associated with states (\ref{colorM}), (\ref{colorB}) and (\ref{colorT}).

The matrices of color factors are available by applying Eq.~(\ref{colorfc}) to wave functions, which give
\begin{equation}\label{cfcM}
    \left\langle \boldsymbol{\lambda_{1}}\cdot \boldsymbol{\lambda_{2}} \right\rangle=-\frac{16}{3}
\end{equation}
for meson in $(\phi^{M})$,
\begin{equation}\label{cfcB}
    \left\langle \boldsymbol{\lambda_{1}}\cdot \boldsymbol{\lambda_{2}} \right\rangle=
    \left\langle \boldsymbol{\lambda_{1}}\cdot \boldsymbol{\lambda_{3}} \right\rangle=
    \left\langle \boldsymbol{\lambda_{2}}\cdot \boldsymbol{\lambda_{3}} \right\rangle=-\frac{8}{3}
\end{equation}
for baryon in $(\phi^{B})$, and
\begin{equation}\label{cfcT}
    \begin{aligned}
        &\left\langle \boldsymbol{\lambda_{1}}\cdot \boldsymbol{\lambda_{2}} \right\rangle=
        \left\langle \boldsymbol{\lambda_{3}}\cdot \boldsymbol{\lambda_{4}} \right\rangle=
        \begin{bmatrix}
            \frac{4}{3} & 0 \\
            0 & -\frac{8}{3}
        \end{bmatrix} \\
        &\left\langle \boldsymbol{\lambda_{1}}\cdot \boldsymbol{\lambda_{3}} \right\rangle=
        \left\langle \boldsymbol{\lambda_{2}}\cdot \boldsymbol{\lambda_{4}} \right\rangle=
        \begin{bmatrix}
            -\frac{10}{3} & 2\sqrt{2} \\
            2\sqrt{2} & -\frac{4}{3}
        \end{bmatrix} \\
        &\left\langle \boldsymbol{\lambda_{1}}\cdot \boldsymbol{\lambda_{4}} \right\rangle=
        \left\langle \boldsymbol{\lambda_{2}}\cdot \boldsymbol{\lambda_{3}} \right\rangle=
        \begin{bmatrix}
            -\frac{10}{3} & -2\sqrt{2} \\
            -2\sqrt{2} & -\frac{4}{3}
        \end{bmatrix}
    \end{aligned}
\end{equation}
for tetraquark in $(\phi_{1}^{T},\phi_{2}^{T})$ space.


\medskip
\section*{Appendix B}\label{AppendixB}
\setcounter{equation}{0}
\renewcommand{\theequation}{B\arabic{equation}}
For meson, baryon, and tetraquark systems labeled $(q_{1}\bar{q}_{2})$ $(q_{1}q_{2})q_{3}$ and $(q_{1}q_{2})(\bar{q}_{3}\bar{q}_{4})$ respectively,
the total spin wave functions collected by checking C-G coefficients read
\begin{equation}\label{pspinM}
    \chi_{1}^{M}=\uparrow\uparrow \quad \chi_{2}^{M}=\frac{1}{\sqrt{2}} \left(\uparrow\downarrow-\downarrow\uparrow\right)
\end{equation}
\begin{equation}\label{pspinB}
    \begin{aligned}
        \chi_{1}^{B}&=\uparrow\uparrow\uparrow \\
        \chi_{2}^{B}&=\sqrt{\frac{2}{3}}\uparrow\uparrow\downarrow-\frac{1}{\sqrt{6}} \left(\uparrow\downarrow\uparrow+\downarrow\uparrow\uparrow\right) \\
        \chi_{3}^{B}&=\frac{1}{\sqrt{2}} \left(\uparrow\downarrow\uparrow-\downarrow\uparrow\uparrow\right)
    \end{aligned}
\end{equation}
\begin{equation}\label{pspinT}
    \begin{aligned}
        \chi_{1}^{T}&=\uparrow\uparrow\uparrow\uparrow \\
        \chi_{2}^{T}&=\frac{1}{2} \left(\uparrow\uparrow\uparrow\downarrow+\uparrow\uparrow\downarrow\uparrow-\uparrow\downarrow\uparrow\uparrow-\downarrow\uparrow\uparrow\uparrow\right) \\
        \chi_{3}^{T}&=\frac{1}{\sqrt{3}} \left(\uparrow\uparrow\downarrow\downarrow+\downarrow\downarrow\uparrow\uparrow\right)\\
        &-\frac{1}{2\sqrt{3}} \left(\uparrow\downarrow\uparrow\downarrow+\uparrow\downarrow\downarrow\uparrow+\downarrow\uparrow\uparrow\downarrow+\downarrow\uparrow\downarrow\uparrow\right) \\
        \chi_{4}^{T}&=\frac{1}{\sqrt{2}} \left(\uparrow\uparrow\uparrow\downarrow-\uparrow\uparrow\downarrow\uparrow\right) \\
        \chi_{5}^{T}&=\frac{1}{\sqrt{2}} \left(\uparrow\downarrow\uparrow\uparrow-\downarrow\uparrow\uparrow\uparrow\right) \\
        \chi_{6}^{T}&=\frac{1}{2} \left(\uparrow\downarrow\uparrow\downarrow-\uparrow\downarrow\downarrow\uparrow-\downarrow\uparrow\uparrow\downarrow+\downarrow\uparrow\downarrow\uparrow\right)
    \end{aligned}
\end{equation}
associated with states (\ref{spinM}), (\ref{spinB}) and (\ref{spinT}).

The matrices of spin factors are available by applying Eq.~(\ref{spinfc}) to wave functions, which give
\begin{equation}\label{sfcM}
    \left\langle \boldsymbol{\sigma_{1}}\cdot \boldsymbol{\sigma_{2}} \right\rangle=
    \begin{bmatrix}
        1 & 0 \\
        0 & -3
    \end{bmatrix}
\end{equation}
for meson in $(\chi_{1}^{M},\chi_{2}^{M})$ space,
    \begin{align}
        \left\langle \boldsymbol{\sigma_{1}}\cdot \boldsymbol{\sigma_{2}} \right\rangle&=
        \begin{bmatrix}
            1 & 0 & 0 \\
            0 & 1 & 0 \\
            0 & 0 & -3
        \end{bmatrix} \\
        \left\langle \boldsymbol{\sigma_{1}}\cdot \boldsymbol{\sigma_{3}} \right\rangle&=
        \begin{bmatrix}
            1 & 0 & 0 \\
            0 & -2 & -\sqrt{3} \\
            0 & -\sqrt{3} & 0
        \end{bmatrix} \\
        \left\langle \boldsymbol{\sigma_{2}}\cdot \boldsymbol{\sigma_{3}} \right\rangle&=
        \begin{bmatrix}
            1 & 0 & 0 \\
            0 & -2 & \sqrt{3} \\
            0 & \sqrt{3} & 0
        \end{bmatrix}
    \end{align}
for baryon in $(\chi_{1}^{B},\chi_{2}^{B},\chi_{3}^{B})$ space, and
    \begin{align}
        \left\langle \boldsymbol{\sigma_{1}}\cdot \boldsymbol{\sigma_{2}} \right\rangle&=
        \begin{bmatrix}
            1 & 0 & 0 & 0 & 0 & 0 \\
            0 & 1 & 0 & 0 & 0 & 0 \\
            0 & 0 & 1 & 0 & 0 & 0 \\
            0 & 0 & 0 & 1 & 0 & 0 \\
            0 & 0 & 0 & 0 & -3 & 0 \\
            0 & 0 & 0 & 0 & 0 & -3
        \end{bmatrix} \\
        \left\langle \boldsymbol{\sigma_{1}}\cdot \boldsymbol{\sigma_{3}} \right\rangle&=
        \begin{bmatrix}
            1 & 0 & 0 & 0 & 0 & 0 \\
            0 & -1 & 0 & \sqrt{2} & -\sqrt{2} & 0 \\
            0 & 0 & -2 & 0 & 0 & -\sqrt{3} \\
            0 & \sqrt{2} & 0 & 0 & 1 & 0 \\
            0 & -\sqrt{2} & 0 & 1 & 0 & 0 \\
            0 & 0 & -\sqrt{3} & 0 & 0 & 0
        \end{bmatrix} \\
        \left\langle \boldsymbol{\sigma_{1}}\cdot \boldsymbol{\sigma_{4}} \right\rangle&=
        \begin{bmatrix}
            1 & 0 & 0 & 0 & 0 & 0 \\
            0 & -1 & 0 & -\sqrt{2} & -\sqrt{2} & 0 \\
            0 & 0 & -2 & 0 & 0 & \sqrt{3} \\
            0 & -\sqrt{2} & 0 & 0 & -1 & 0 \\
            0 & -\sqrt{2} & 0 & -1 & 0 & 0 \\
            0 & 0 & \sqrt{3} & 0 & 0 & 0
        \end{bmatrix} \\
        \left\langle \boldsymbol{\sigma_{2}}\cdot \boldsymbol{\sigma_{3}} \right\rangle&=
        \begin{bmatrix}
            1 & 0 & 0 & 0 & 0 & 0 \\
            0 & -1 & 0 & \sqrt{2} & \sqrt{2} & 0 \\
            0 & 0 & -2 & 0 & 0 & \sqrt{3} \\
            0 & \sqrt{2} & 0 & 0 & -1 & 0 \\
            0 & \sqrt{2} & 0 & -1 & 0 & 0 \\
            0 & 0 & \sqrt{3} & 0 & 0 & 0
        \end{bmatrix} \\
        \left\langle \boldsymbol{\sigma_{2}}\cdot \boldsymbol{\sigma_{4}} \right\rangle&=
        \begin{bmatrix}
            1 & 0 & 0 & 0 & 0 & 0 \\
            0 & -1 & 0 & -\sqrt{2} & \sqrt{2} & 0 \\
            0 & 0 & -2 & 0 & 0 & -\sqrt{3} \\
            0 & -\sqrt{2} & 0 & 0 & 1 & 0 \\
            0 & \sqrt{2} & 0 & 1 & 0 & 0 \\
            0 & 0 & -\sqrt{3} & 0 & 0 & 0
        \end{bmatrix} \\
        \left\langle \boldsymbol{\sigma_{3}}\cdot \boldsymbol{\sigma_{4}} \right\rangle&=
        \begin{bmatrix}
            1 & 0 & 0 & 0 & 0 & 0 \\
            0 & 1 & 0 & 0 & 0 & 0 \\
            0 & 0 & 1 & 0 & 0 & 0 \\
            0 & 0 & 0 & -3 & 0 & 0 \\
            0 & 0 & 0 & 0 & 1 & 0 \\
            0 & 0 & 0 & 0 & 0 & -3
        \end{bmatrix} \label{sfcT}
    \end{align}
for tetraquark in $(\chi_{1}^{T},\chi_{2}^{T},\chi_{3}^{T},\chi_{4}^{T},\chi_{5}^{T},\chi_{6}^{T})$ space.

For now, the CMI for a hadronic state can be expressed by summing up all the $C_{ij}$ by color and spin factors,
which are diagonal elements of matrices in Eq.~(\ref{cfcM}-\ref{cfcT}) and Eq.~(\ref{sfcM}-\ref{sfcT}).
The off-diagonal elements are used for chromomagnetic mixing.


\medskip
\section*{Appendix C}\label{AppendixC}
\setcounter{equation}{0}
\renewcommand{\theequation}{C\arabic{equation}}
According to Appendix A and B, we list CMI matrices of all the chromomagnetic mixing involved in this paper, such as
\begin{equation}\label{mixing23}
    \begin{pmatrix}
        \frac{8}{3}C_{12}-\frac{16}{3}C_{13}-\frac{16}{3}C_{23} & -\frac{8\sqrt{3}}{3}C_{13}+\frac{8\sqrt{3}}{3}C_{23} \\
        -\frac{8\sqrt{3}}{3}C_{13}+\frac{8\sqrt{3}}{3}C_{23} & -8C_{12} \\
    \end{pmatrix}
\end{equation}
for $(\phi^{B}\chi_{2}^{B},\phi^{B}\chi_{3}^{B})$ mixing,
\begin{equation}\label{21mixing36}
    \begin{pmatrix}
        \frac{8}{3}(\alpha-\beta)   & 2\sqrt{6}\beta \\
        2\sqrt{6}\beta              & 4\alpha
    \end{pmatrix}
\end{equation}
for $(\phi_{2}^{T}\chi_{3}^{T},\phi_{1}^{T}\chi_{6}^{T})$ mixing,
\begin{equation}\label{21mixing54}
    \begin{pmatrix}
        -\frac{8}{3}\theta  & -2\sqrt{2}\beta \\
        -2\sqrt{2}\beta     & -\frac{4}{3}\eta
    \end{pmatrix}
\end{equation}
for $(\phi_{2}^{T}\chi_{5}^{T},\phi_{1}^{T}\chi_{4}^{T})$ mixing,
\begin{equation}\label{12mixing36}
    \begin{pmatrix}
        -\frac{4}{3}(\alpha+5\beta) & 2\sqrt{6}\beta \\
        2\sqrt{6}\beta              & -8\alpha
    \end{pmatrix}
\end{equation}
for $(\phi_{1}^{T}\chi_{3}^{T},\phi_{2}^{T}\chi_{6}^{T})$ mixing,
\begin{equation}\label{221mixing245}
    \begin{pmatrix}
        \frac{4}{3}(2\alpha-\beta)  & \frac{4\sqrt{2}}{3}\delta & 4\delta \\
        \frac{4\sqrt{2}}{3}\delta   & \frac{8}{3}\eta           & -2\sqrt{2}\beta \\
        4\delta                     & -2\sqrt{2}\beta           & \frac{4}{3}\beta
    \end{pmatrix}
\end{equation}
for $(\phi_{2}^{T}\chi_{2}^{T},\phi_{2}^{T}\chi_{4}^{T},\phi_{1}^{T}\chi_{5}^{T})$ mixing,
\begin{equation}\label{221mixing254}
    \begin{pmatrix}
        \frac{4}{3}(2\alpha-\beta)  & -\frac{4\sqrt{2}}{3}\gamma    & -4\gamma \\
        -\frac{4\sqrt{2}}{3}\gamma  & -\frac{8}{3}\theta            & -2\sqrt{2}\beta \\
        -4\gamma                    & -2\sqrt{2}\beta               & -\frac{4}{3}\eta
    \end{pmatrix}
\end{equation}
for $(\phi_{2}^{T}\chi_{2}^{T},\phi_{2}^{T}\chi_{5}^{T},\phi_{1}^{T}\chi_{1}^{T})$ mixing,
\begin{equation}\label{112mixing245}
    \begin{pmatrix}
        -\frac{2}{3}(2\alpha+5\beta)    & \frac{10\sqrt{2}}{3}\delta    & 4\delta \\
        \frac{10\sqrt{2}}{3}\delta      & -\frac{4}{3}\eta              & -2\sqrt{2}\beta \\
        4\delta                         & -2\sqrt{2}\beta               & -\frac{8}{3}\theta
    \end{pmatrix}
\end{equation}
for $(\phi_{1}^{T}\chi_{2}^{T},\phi_{1}^{T}\chi_{4}^{T},\phi_{2}^{T}\chi_{5}^{T})$ mixing,
where $\alpha=C_{12}+C_{34}$, $\beta=C_{13}+C_{14}+C_{23}+C_{24}$, $\gamma=C_{13}+C_{14}-C_{23}-C_{24}$, $\delta=C_{13}-C_{14}+C_{23}-C_{24}$,
$\eta=C_{12}-3C_{34}$ and $\theta=3C_{12}-C_{34}$, refering Ref.~\cite{Luo:2017eub}.

And we also present binding energies
\begin{equation}
    B_{12}+B_{13}+B_{23}
\end{equation}
for $\phi^{B}$,
\begin{equation}
    -\frac{1}{2}B_{12}+\frac{5}{4}B_{13}+\frac{5}{4}B_{14}+\frac{5}{4}B_{23}+\frac{5}{4}B_{24}-\frac{1}{2}B_{34}
\end{equation}
for $\phi_{1}^{T}$, and
\begin{equation}
    B_{12}+\frac{1}{2}B_{13}+\frac{1}{2}B_{14}+\frac{1}{2}B_{23}+\frac{1}{2}B_{24}+B_{34}
\end{equation}
for $\phi_{2}^{T}$ respectively, where $B_{ij}$ takes value from Eq.~(\ref{bingdings}) and equals 0 if there is no short-distance binding.


\begin{thebibliography}{99}

\bibitem{Aaij:2017ueg}
    R.~Aaij \textit{et al.} [LHCb],
    %``Observation of the doubly charmed baryon $\Xi_{cc}^{++}$,''
    Phys. Rev. Lett. \textbf{119}, no.11, 112001 (2017)
    doi:10.1103/PhysRevLett.119.112001
    [arXiv:1707.01621 [hep-ex]].

\bibitem{KR:prl17} M. Karliner and J. L. Rosner, Phys.Rev.Lett.\textbf{119}, 202001(2017)
\bibitem{EichtenQ:prl17} Estia J. Eichten and C. Quigg, Phys. Rev. Lett.\textbf{119}, 202002(2017)
\bibitem{LuoCL:Epj17} S. Q. Lou, K. Chen, X. Liu, Y. R. Liu and S. L. Zhi, Eur. Phys. J. C\textbf{77}, 709(2017),

\bibitem{DeGrandJ:1975D}
    T. DeGrand, R. L. Jaffe et al., 
    Phys. Rev. D \textbf{12}, 2060 (1975).

\bibitem{DeGrand:1975cf}
    T.~A.~DeGrand, R.~L.~Jaffe, K.~Johnson and J.~E.~Kiskis,
    %``Masses and Other Parameters of the Light Hadrons,''
    Phys. Rev. D \textbf{12}, 2060 (1975)
    doi:10.1103/PhysRevD.12.2060

\bibitem{Strottman:1979qu}
    D.~Strottman,
    %``Multi - Quark Baryons and the MIT Bag Model,''
    Phys. Rev. D \textbf{20}, 748-767 (1979)
    doi:10.1103/PhysRevD.20.748
    
\bibitem{FabbroJR:D05}    
    A. Del Fabbro, D. Janc, M. Rosina, D. Treleani, Phys. Rev. D 71 (2005) 014008, arXiv:hep-ph/0408258.

\bibitem{Carlson:1982er}
    C.~E.~Carlson, T.~H.~Hansson and C.~Peterson,
    %``Meson, Baryon and Glueball Masses in the MIT Bag Model,''
    Phys. Rev. D \textbf{27}, 1556-1564 (1983)
    doi:10.1103/PhysRevD.27.1556

\bibitem{Haxton:1980mc}
    W.~C.~Haxton and L.~Heller,
    %``The Heavy Quark - Anti-quark Potential in the {MIT} Bag Model,''
    Phys. Rev. D \textbf{22}, 1198 (1980)
    doi:10.1103/PhysRevD.22.1198

\bibitem{Aerts:1980rf}
    A.~T.~M.~Aerts and L.~Heller,
    %``The Potential Energy of Three Heavy Quarks in the {MIT} Bag Model,''
    Phys. Rev. D \textbf{23}, 185 (1981)
    doi:10.1103/PhysRevD.23.185

\bibitem{Fleck:1989mb}
    S.~Fleck and J.~M.~Richard,
    %``Baryons with double charm,''
    Prog. Theor. Phys. \textbf{82}, 760-774 (1989)
    doi:10.1143/PTP.82.760

\bibitem{He:2004px}
    D.~H.~He, K.~Qian, Y.~B.~Ding, X.~Q.~Li and P.~N.~Shen,
    %``Evaluation of spectra of baryons containing two heavy quarks in bag model,''
    Phys. Rev. D \textbf{70}, 094004 (2004)
    doi:10.1103/PhysRevD.70.094004
    [arXiv:hep-ph/0403301 [hep-ph]].

\bibitem{Bernotas:2012nz}
    A.~Bernotas and V.~Simonis,
    %``Magnetic moments of heavy baryons in the bag model reexamined,''
    doi:10.3952/physics.v53i2.2668
    [arXiv:1209.2900 [hep-ph]].
 \bibitem{Meng:201HHH:2021}   
    Q. Meng,1 M. Harada,2, 3, 4 E. Hiyama,5, 6, 4, 7 A. Hosaka,7, 4, 6 and M. Oka, Doubly Heavy Tetraquark Resonant States, arXiv:2106.11868[hep-ph]

\bibitem{Moinester:2002uw}
    M.~A.~Moinester \textit{et al.} [SELEX],
    %``First Observation of Doubly Charmed Baryons,''
    Czech. J. Phys. \textbf{53}, B201-B213 (2003)
    [arXiv:hep-ex/0212029 [hep-ex]].

\bibitem{Ebert:2004ck}
    D.~Ebert, R.~N.~Faustov, V.~O.~Galkin and A.~P.~Martynenko,
    %``Semileptonic decays of doubly heavy baryons in the relativistic quark model,''
    Phys. Rev. D \textbf{70}, 014018 (2004)
    [erratum: Phys. Rev. D \textbf{77}, 079903 (2008)]
    doi:10.1103/PhysRevD.70.014018
    [arXiv:hep-ph/0404280 [hep-ph]].

\bibitem{Karliner:2014gca}
    M.~Karliner and J.~L.~Rosner,
    %``Baryons with two heavy quarks: Masses, production, decays, and detection,''
    Phys. Rev. D \textbf{90}, no.9, 094007 (2014)
    doi:10.1103/PhysRevD.90.094007
    [arXiv:1408.5877 [hep-ph]].

\bibitem{Liu:2018euh}
    M.~Z.~Liu, Y.~Xiao and L.~S.~Geng,
    %``Magnetic moments of the spin-1/2 doubly charmed baryons in covariant baryon chiral perturbation theory,''
    Phys. Rev. D \textbf{98}, no.1, 014040 (2018)
    doi:10.1103/PhysRevD.98.014040
    [arXiv:1807.00912 [hep-ph]].

\bibitem{Luo:2017eub}
    S.~Q.~Luo, K.~Chen, X.~Liu, Y.~R.~Liu and S.~L.~Zhu,
    %``Exotic tetraquark states with the $qq\bar{Q}\bar{Q}$ configuration,''
    Eur. Phys. J. C \textbf{77}, no.10, 709 (2017)
    doi:10.1140/epjc/s10052-017-5297-4
    [arXiv:1707.01180 [hep-ph]].
\bibitem{AliLS:ppn2017}    
    A. Ali, J. S. Lange, S. Stone, Exotics: Heavy pentaquarks and tetraquarks, Prog. Part. Nucl. Phys. \textbf{97} (2017) 123�C198
    
\bibitem{Liu:2019zoy}
    Y.~R.~Liu, H.~X.~Chen, W.~Chen, X.~Liu and S.~L.~Zhu,
    %``Pentaquark and Tetraquark states,''
    Prog. Part. Nucl. Phys. \textbf{107}, 237-320 (2019)
    doi:10.1016/j.ppnp.2019.04.003
    [arXiv:1903.11976 [hep-ph]].

\bibitem{Chodos:1974pn}
    A.~Chodos, R.~L.~Jaffe, K.~Johnson and C.~B.~Thorn,
    %``Baryon Structure in the Bag Theory,''
    Phys. Rev. D \textbf{10}, 2599 (1974)
    doi:10.1103/PhysRevD.10.2599

\bibitem{Karliner:17N}
    M.~Karliner and J.~L.~Rosner, Quark-level analogue of nuclear fusion with doubly heavy baryons,
    Nature \textbf{551}, 89 (2017)
    doi:10.1038/nature24289
    arXiv:1708.02547 [hep-ph]    
    
\bibitem{Karliner:18D}
    M.~Karliner and J.~L.~Rosner, Scaling of P-wave excitation energies in heavy-quark systems, 
    Phys. Rev. D \textbf{98}, no.9, 074026 (2018)
    doi:10.1103/PhysRevD.98.074026
    arXiv:1808.07869 [hep-ph] 

\bibitem{Wang:2016dzu}
    G.~J.~Wang, R.~Chen, L.~Ma, X.~Liu and S.~L.~Zhu,
    %``Magnetic moments of the hidden-charm pentaquark states,''
    Phys. Rev. D \textbf{94}, no.9, 094018 (2016)
    doi:10.1103/PhysRevD.94.094018
    [arXiv:1605.01337 [hep-ph]].

\bibitem{Ebert:2002pp}
    D.~Ebert, R.~N.~Faustov and V.~O.~Galkin,
    %``Properties of heavy quarkonia and $B_c$ mesons in the relativistic quark model,''
    Phys. Rev. D \textbf{67}, 014027 (2003)
    doi:10.1103/PhysRevD.67.014027
    [arXiv:hep-ph/0210381 [hep-ph]].

\bibitem{LHCb:2020kd} LHC Seminar,
    $B\to D\bar{D}h$ decays: ''A new (virtual) laboratory for exotic particle searches at LHCb'',
    by Daniel Johnson, CERN, August 11, 2020,
    \href{https://indico.cern.ch/event/900975/}{https://indico.cern.ch/event/900975/}.

\bibitem{Karliner:2020vsi}
    M.~Karliner and J.~L.~Rosner,
    %``First exotic hadron with open heavy flavor: $cs\bar u\bar d$ tetraquark,''
    Phys. Rev. D \textbf{102}, no.9, 094016 (2020)
    doi:10.1103/PhysRevD.102.094016
    [arXiv:2008.05993 [hep-ph]].

\bibitem{Lu:2020rog}
    Q.~F.~L\"u, D.~Y.~Chen and Y.~B.~Dong,
    %``Masses of doubly heavy tetraquarks $T_{QQ^\prime}$ in a relativized quark model,''
    Phys. Rev. D \textbf{102}, no.3, 034012 (2020)
    doi:10.1103/PhysRevD.102.034012
    [arXiv:2006.08087 [hep-ph]].

\bibitem{Ebert:2007rn}
    D.~Ebert, R.~N.~Faustov, V.~O.~Galkin and W.~Lucha,
    %``Masses of tetraquarks with two heavy quarks in the relativistic quark model,''
    Phys. Rev. D \textbf{76}, 114015 (2007)
    doi:10.1103/PhysRevD.76.114015
    [arXiv:0706.3853 [hep-ph]].

\end{thebibliography}

\end{document}
